\documentclass[11pt,reqno]{amsart}

\usepackage[T1]{fontenc}
\usepackage[utf8]{inputenc}
\usepackage{amsmath}
\usepackage{amssymb}
\usepackage{amsthm}
\usepackage{mathtools}
\usepackage[margin=1.15in]{geometry}
\usepackage[colorlinks=true,linkcolor=blue,citecolor=blue,urlcolor=blue]{hyperref}
\usepackage[capitalise,nameinlink]{cleveref}

%---------------------------------------------------------------
% Theorem environments
%---------------------------------------------------------------
\theoremstyle{plain}
\newtheorem{theorem}{Theorem}[section]
\newtheorem{proposition}[theorem]{Proposition}
\newtheorem{lemma}[theorem]{Lemma}
\newtheorem{corollary}[theorem]{Corollary}

\theoremstyle{definition}
\newtheorem{definition}[theorem]{Definition}
\newtheorem{example}[theorem]{Example}
\newtheorem{remark}[theorem]{Remark}

\crefname{theorem}{Theorem}{Theorems}
\crefname{proposition}{Proposition}{Propositions}
\crefname{lemma}{Lemma}{Lemmas}
\crefname{corollary}{Corollary}{Corollaries}
\crefname{definition}{Definition}{Definitions}
\crefname{example}{Example}{Examples}
\crefname{remark}{Remark}{Remarks}
\crefname{section}{Section}{Sections}
\crefname{equation}{Equation}{Equations}

%---------------------------------------------------------------
% Macros
%---------------------------------------------------------------
\newcommand{\Rr}{\mathbb{R}}
\newcommand{\Cc}{\mathbb{C}}
\newcommand{\Zz}{\mathbb{Z}}
\newcommand{\Uone}{\mathrm{U}(1)}
\newcommand{\Src}{S}
\newcommand{\Rdt}{T}
\newcommand{\Pow}{\mathfrak{P}}          % power set -- deliberately distinct from probability notation
\newcommand{\CompatSet}{\mathsf{A}}      % compatible-answer set
\newcommand{\Spec}{\operatorname{Spec}}
\newcommand{\Img}{\operatorname{Im}}
\newcommand{\id}{\operatorname{id}}
\newcommand{\Dx}{D_{x}}
\DeclareMathOperator{\Hol}{Hol}

% Allow line breaks inside recurring hyphenated compounds
\hyphenation{tar-get rel-a-tiv-i-ty read-out sol-va-bil-i-ty
  iso-spec-tral hol-o-nom-y cur-va-ture}

%---------------------------------------------------------------
\begin{document}

\title[Target-Relative Necessity of Completion]{Target-Relative Necessity of Completion:\\
When Readout Loss Obstructs, and What a Sufficient\\ Extension Must Retain}

\author{Jeremy Rodgers}
\thanks{Website: \href{https://everythingequation.com}{everythingequation.com}}
\address{Independent Researcher}

\date{June 2026}

\subjclass[2020]{Primary 58J50, 81Q70; Secondary 58J53, 81Q10, 53C29}

\keywords{Quotient readout, information loss, target solvability, flat connections, holonomy,
covariant Laplacian, spectral rigidity, projected dynamics}

\begin{abstract}
A surjective, noninjective readout map $p\colon\Src\to\Rdt$ from a reduced source space
necessarily discards distinctions among reduced source states. It does not follow that the
discarded distinctions matter. Whether they matter is a property of the question being asked,
not of the readout alone. This paper makes that dependence precise. A \emph{target} is specified
by a correct-answer map $a_Q\colon\Src\to A_Q$, and we prove that $Q$ admits an exact
readout-only solution if and only if $a_Q$ is constant on every fiber of $p$; consequently a
lost source relation obstructs $Q$ precisely when it is \emph{active} for $Q$, i.e.\ when the
realized answer varies across some readout fiber. When exactness fails, the compatible-answer
set $\CompatSet_Q(t)=a_Q(p^{-1}(t))$ is shown to be the unique pointwise-least uniformly sound
set-valued rule available without further information; it is a set-theoretic, not probabilistic,
object. We then characterize what any repair must contain. Every surjective readout extension
$e\colon\Src\to E$ through which $a_Q$ factors must separate every pair of source states with
different correct answers, and the joint target image $E_Q:=\Img(p,a_Q)$ is itself such an
extension and is universally coarsest: every $Q$-sufficient extension $e$ admits a unique
surjection $\varphi_E\colon E\to E_Q$ with $j_Q=\varphi_E\circ e$. The framework is realized
concretely for flat $\Uone$ connections on a circle. Curvature is a deliberately restricted
readout that collapses the entire flat moduli space $\Src_{\mathrm{flat}}\cong\Uone$ to a point,
yet the spectrum of the charged covariant Laplacian $H_\theta=-\Dx^2$ separates $\theta=0$ from
$\theta=\pi$; we prove that $\Spec(H_\theta)=\Spec(H_{\theta'})$, with multiplicities, if and
only if $\theta'\equiv\pm\theta \pmod{2\pi}$, so that the coarsest completion sufficient for the
spectral target is the inversion quotient coordinatized by $\cos\theta$, strictly coarser than
full holonomy. Finally the same criterion is applied to finite-time projected dynamics: an
autonomous readout propagator $U_\tau$ with $p\circ\Phi_\tau=U_\tau\circ p$ exists if and only
if the future readout is constant on present readout fibers.
\end{abstract}

\maketitle

%===============================================================
\section{Introduction}\label{sec:intro}

\subsection{The question left open by exactness}

Physical descriptions are routinely obtained by passing from a rich source description to a
coarser one. Gauge fixing, symmetry reduction, coarse graining, restriction to a chosen family
of observables, and projection onto a subsystem all have the same formal shape: a surjection
\begin{equation}\label{eq:readout-intro}
p\colon\Src\longrightarrow\Rdt
\end{equation}
from a space $\Src$ of admissible states onto a space $\Rdt$ of readout values. The companion
paper to the present one, on source--readout non-equivalence \cite{ShadowI}, established two
facts about such maps. First, an exact quotient readout need not be equivalent to the reduced
source it was computed from: even when no approximation is committed and no information is
thrown away by error, the map $p$ may fail to be injective, and then $\Rdt$ simply does not
carry all the distinctions that $\Src$ carries. Second, a source-defined relation
$r\colon\Src\to V$ descends to a relation on $\Rdt$ precisely when $r$ is constant on the fibers
$F_t:=p^{-1}(t)$.

That is a clean criterion, but by itself it is a statement about the relation $r$, and it is
tempting to over-read it. The over-reading runs: since $p$ is noninjective, information is lost;
since information is lost, the readout description is defective; since the readout description
is defective, it must be repaired by additional variables. Each step of this chain is either
false or unsupported. Noninjectivity is a structural fact about $p$, and structural loss is
cheap: almost every useful description loses something. The relevant question is never whether
$p$ loses distinctions, but whether it loses \emph{the distinctions that the question at hand
depends on}. A thermodynamic readout that forgets molecular positions is not thereby unfit to
answer thermodynamic questions.

The present paper is devoted to making that dependence exact. We ask:

\begin{quote}
Given a nominated mathematical or physical question, when does the information discarded by the
readout genuinely obstruct answering it, and, when it does, what must \emph{any} successful
repair retain?
\end{quote}

Both halves of this question are answered here by positive mathematics. In particular the
impossibility half is not obtained by assuming that no readout-only construction exists. It is
\emph{derived} from the observation that a deterministic construction taking only the readout as
input is necessarily of the form $b\circ p$, hence constant on fibers, whereas the correct answer
may vary along a fiber. Variation of the correct-answer map on a single fiber is the entire
source of the obstruction, and the obstruction is therefore exactly as strong, and no stronger,
than that variation.

\subsection{Targets and correct-answer maps}

The technical device that makes the theory target-relative is elementary and, we will argue,
unavoidable. A \emph{target} $Q$ is not a piece of prose. It is a pair, consisting of a set
$A_Q$ of admissible answers together with a map
\begin{equation}\label{eq:answermap-intro}
a_Q\colon\Src\longrightarrow A_Q
\end{equation}
assigning to each reduced source state the answer that is in fact correct for that state. Nothing
mysterious is intended by ``correct'': if the target is ``does the charged Laplacian have a zero
mode?'' then $A_Q=\{\text{yes},\text{no}\}$ and $a_Q$ records the true answer for each
connection; if the target is ``what is the spectrum, with multiplicities?'' then $A_Q$ is a set
of multisets and $a_Q$ records the actual spectrum. The values of $a_Q$ may be Boolean,
numerical, functional, geometric, spectral, or dynamical; the theory is indifferent.

What the device buys is precision about the object of the obstruction. ``Information is lost''
is a statement about $p$. ``$Q$ cannot be solved from the readout'' is a statement about the
pair $(p,a_Q)$, and \cref{thm:exact-solvability} shows the second is equivalent to the failure
of $a_Q$ to be fiber-constant. Two different targets, evaluated against the same lossy readout,
may receive opposite verdicts, and in \cref{sec:flat,sec:spectrum,sec:minimal-spectral} we
exhibit exactly that: a curvature readout that is perfectly adequate for one class of questions
about flat $\Uone$ connections and provably inadequate for another, together with a completion
that is minimal for the second class and strictly insufficient for a third.

\subsection{Contributions}

The paper is organized around a single logical progression, and we summarize it here so that the
dependency order is visible from the outset.

\begin{enumerate}
\item \Cref{sec:setup} fixes the reduced source space, the readout, and the target, and insists
that all statements concern reduced source states --- states already quotiented by every declared
gauge and coordinate redundancy --- so that no apparent ``loss'' is an artifact of retaining
redundant labels.

\item \Cref{sec:exact} proves the exact solvability criterion
(\cref{thm:exact-solvability}): $Q$ is exactly solvable from the readout iff $a_Q$ is constant
on every fiber of $p$; and derives the obstruction (\cref{cor:obstruction}) in the form actually
needed, namely that a single fiber on which $a_Q$ takes two values already defeats every
deterministic readout-only construction. We show additionally
(\cref{prop:auxiliary-useless}) that auxiliary quantities computed from the readout cannot repair
this, because they too factor through $p$.

\item \Cref{sec:active} separates \emph{structural} loss from \emph{target-active} loss. A
source relation $r$ is active for $Q$ on a fiber when the realized answers it induces there
differ; it is inactive when they agree, \emph{even if $r$ itself varies}. Activity on one fiber
obstructs (\cref{thm:active-obstructs}); global inactivity preserves exact solvability
(\cref{thm:inactive-preserves}).

\item \Cref{sec:uncertainty} identifies what is still available when exactness fails: the
compatible-answer set $\CompatSet_Q(t)=a_Q(F_t)$, proved to be uniformly sound and pointwise
contained in every uniformly sound set-valued rule (\cref{thm:minimal-sound}). We stress that
$\Pow(A_Q)$ is the power set and that this is a set-valued, not probabilistic, statement.

\item \Cref{sec:extensions,sec:separation} define readout extensions and target sufficiency.
We then prove the necessary-separation theorem (\cref{thm:separation}): every $Q$-sufficient
extension must separate each pair of source states whose correct answers differ.

\item \Cref{sec:joint} constructs the joint target image $E_Q=\Img(p,a_Q)$, proves it is a
$Q$-sufficient extension, and proves the universal factorization
(\cref{thm:universal}) in the correct direction: every $Q$-sufficient extension $E$ maps onto
$E_Q$, not conversely. $E_Q$ is therefore the coarsest sufficient completion for the single
target $Q$.

\item \Cref{sec:flat,sec:spectrum,sec:minimal-spectral} realize all of this for flat $\Uone$
connections on $S^1$ with a curvature readout, a charged covariant Laplacian, and its spectrum,
culminating in the spectral rigidity result $\Spec(H_\theta)=\Spec(H_{\theta'})
\iff \theta'\equiv\pm\theta$.

\item \Cref{sec:dynamics} applies the same criterion to autonomous projected dynamics at finite
time.
\end{enumerate}

\Cref{sec:discussion} delimits scope. The sequel \cite{ShadowIII} generalizes from a single
correct-answer map to families of realization relations; continuous-time and memory-kernel
questions belong to \cite{ShadowV}. Neither construction is used here, and the results below are
self-contained.

%===============================================================
\section{Reduced source, readout, and target}\label{sec:setup}

\subsection{The reduced source space}

Throughout, $\Src$ denotes a \emph{reduced} source space. By this we mean the space obtained
from whatever raw description one began with \emph{after} quotienting by every redundancy that
has been declared physically inert: local gauge transformations, coordinate reparametrizations,
relabelings of identical constituents, and any further identification that the modeling
commitments of the problem require. The point of insisting on this is not fastidiousness. If
$\Src$ still carried gauge copies, then $p$ would appear to lose information whenever it merely
forgot a gauge label, and every theorem below would be vacuous or misleading: one would be
proving obstructions that could be dissolved by a change of gauge. By working on the reduced
space we guarantee that two distinct points of $\Src$ are genuinely distinct physical
configurations, so that a fiber of $p$ containing two points is a genuine collapse of physical
distinctions and not a bookkeeping artifact.

\begin{definition}[Readout]\label{def:readout}
A \emph{readout} is a surjective map $p\colon\Src\to\Rdt$ of sets. For $t\in\Rdt$ we write
\begin{equation}\label{eq:fiber}
F_t := p^{-1}(t) = \{\, s\in\Src : p(s)=t \,\},
\end{equation}
and call $F_t$ the \emph{readout fiber} over $t$. Surjectivity of $p$ means every fiber is
nonempty, and the fibers partition $\Src$.
\end{definition}

We write $\sim_p$ for the equivalence relation on $\Src$ induced by $p$, so $s\sim_p s'$ iff
$p(s)=p(s')$ iff $s$ and $s'$ lie in a common fiber. The readout is \emph{lossy}, or
noninjective, when some fiber contains more than one point. No topological, smooth, or measurable
structure is required for the results of
\cref{sec:exact,sec:active,sec:uncertainty,sec:extensions,sec:separation,sec:joint}; those are
statements about maps of sets, and their strength derives from that generality. Structure is
imposed only in \cref{sec:flat,sec:spectrum,sec:minimal-spectral}, where the source space is a
moduli space of connections and the target values are spectra.

\subsection{Targets}

\begin{definition}[Target and correct-answer map]\label{def:target}
A \emph{target} $Q$ on $\Src$ consists of
\begin{enumerate}
\item a set $A_Q$ of \emph{admissible answers}, and
\item a map $a_Q\colon\Src\to A_Q$, the \emph{correct-answer map}, assigning to each reduced
source state $s$ the answer $a_Q(s)$ that is correct for $s$.
\end{enumerate}
\end{definition}

That $a_Q$ is a well-defined map on all of $\Src$ is a substantive standing assumption, and it is
worth saying what it does and does not import. It asserts that the question has a determinate
answer for each reduced source state --- that once the physical configuration is fixed, so is the
correct answer. It does \emph{not} assert that the answer is computable, that it is known, or
that it is accessible from any particular data. Indeed, the entire content of the paper is the
analysis of when the answer, though determinate as a function on $\Src$, fails to be a function
of $p(s)$. If a purported target does not have a determinate answer per reduced source state ---
if, say, the answer depends on further facts not encoded in $\Src$ --- then it is not a target on
$\Src$ in the sense of \cref{def:target}, and the source space must be enlarged before the
present analysis applies.

\begin{definition}[Exact readout solvability]\label{def:solvable}
The target $Q$ is \emph{exactly solvable from the readout} $p$ if there exists a map
\begin{equation}\label{eq:solvable}
\bar a_Q\colon\Rdt\longrightarrow A_Q
\qquad\text{such that}\qquad
a_Q=\bar a_Q\circ p .
\end{equation}
We then say that $a_Q$ \emph{descends} along $p$, and call $\bar a_Q$ a \emph{readout-only
solution} of $Q$.
\end{definition}

\begin{remark}[Solving a target is weaker than reconstructing the source]\label{rem:weaker}
\Cref{def:solvable} asks only that the correct answer be recoverable from $t=p(s)$. It does not
ask that $s$ itself be recoverable. This is the asymmetry that makes the whole subject
nontrivial. A readout can be violently noninjective --- collapsing an infinite-dimensional fiber
to a point --- and still solve a target exactly, provided the target is blind to everything the
fiber varies over. Conversely, an almost-injective readout that fails to separate a single
critical pair of states can fail to solve a target that hinges on that pair. Injectivity of $p$
is therefore sufficient for exact solvability of \emph{every} target, but it is wildly stronger
than necessary for any \emph{given} target, and the gap between the two is where the useful
mathematics lives.
\end{remark}

%===============================================================
\section{Exact target solvability and its obstruction}\label{sec:exact}

\subsection{The criterion}

\begin{theorem}[Exact target-solvability criterion]\label{thm:exact-solvability}
Let $p\colon\Src\to\Rdt$ be a readout and let $Q=(A_Q,a_Q)$ be a target on $\Src$. The following
are equivalent.
\begin{enumerate}
\item\label{it:solv} $Q$ is exactly solvable from the readout; that is, there is
$\bar a_Q\colon\Rdt\to A_Q$ with $a_Q=\bar a_Q\circ p$.
\item\label{it:fibconst} $a_Q$ is constant on every fiber of $p$: for all $t\in\Rdt$ and all
$s,s'\in F_t$, $a_Q(s)=a_Q(s')$.
\item\label{it:imp} For all $s,s'\in\Src$, $\;p(s)=p(s') \implies a_Q(s)=a_Q(s')$.
\end{enumerate}
Moreover, when these hold, $\bar a_Q$ is unique.
\end{theorem}

\begin{proof}
Conditions \eqref{it:fibconst} and \eqref{it:imp} are restatements of one another: $s,s'$ lie in
a common fiber precisely when $p(s)=p(s')$.

\eqref{it:solv}$\Rightarrow$\eqref{it:imp}. Suppose $a_Q=\bar a_Q\circ p$ and $p(s)=p(s')$.
Applying $\bar a_Q$ to both sides of $p(s)=p(s')$ gives
$a_Q(s)=\bar a_Q(p(s))=\bar a_Q(p(s'))=a_Q(s')$.

\eqref{it:imp}$\Rightarrow$\eqref{it:solv}. We construct $\bar a_Q$. Let $t\in\Rdt$. Since $p$ is
surjective, $F_t\neq\emptyset$; choose $s\in F_t$ and set
\begin{equation}\label{eq:barconstruct}
\bar a_Q(t) := a_Q(s).
\end{equation}
This is independent of the choice of $s$: if $s'\in F_t$ is another element, then $p(s)=t=p(s')$,
so by \eqref{it:imp} $a_Q(s)=a_Q(s')$. Hence $\bar a_Q\colon\Rdt\to A_Q$ is a well-defined map.
(That some element may be chosen from each nonempty fiber is all that is used; no choice
principle beyond the definition of a well-defined value is required, since the value assigned is
independent of the representative.) Now for any $s\in\Src$, put $t=p(s)$; then $s\in F_t$, and by
\eqref{eq:barconstruct} together with representative-independence, $\bar a_Q(p(s))=a_Q(s)$. Thus
$a_Q=\bar a_Q\circ p$.

Uniqueness. Suppose $\bar a_Q$ and $\bar a_Q'$ both satisfy \eqref{eq:solvable}. Let $t\in\Rdt$
and choose $s\in F_t$, possible by surjectivity. Then
$\bar a_Q(t)=\bar a_Q(p(s))=a_Q(s)=\bar a_Q'(p(s))=\bar a_Q'(t)$. As $t$ was arbitrary,
$\bar a_Q=\bar a_Q'$.
\end{proof}

\Cref{thm:exact-solvability} is the fiber--descent theorem of \cite{ShadowI} specialized to the
correct-answer map, and we have given the argument in full so that the present paper stands on
its own. What is new here is not the criterion but the object it is applied to: by taking the
descending map to be the \emph{answer} to a nominated question rather than an arbitrary source
relation, the criterion becomes a statement about the adequacy of the readout \emph{for that
question}, and hence acquires the target-relativity that the rest of the paper develops.

\subsection{The obstruction}

\begin{corollary}[Exact target obstruction]\label{cor:obstruction}
Suppose there exist $s_0,s_1\in\Src$ with
\begin{equation}\label{eq:obstruction-hyp}
p(s_0)=p(s_1)
\qquad\text{and}\qquad
a_Q(s_0)\neq a_Q(s_1).
\end{equation}
Then $Q$ is not exactly solvable from the readout. Explicitly, for every map
$b\colon\Rdt\to A_Q$ there exists $i\in\{0,1\}$ with $b(p(s_i))\neq a_Q(s_i)$: no deterministic
readout-only construction is correct at both $s_0$ and $s_1$.
\end{corollary}

\begin{proof}
Let $b\colon\Rdt\to A_Q$ be arbitrary and write $t:=p(s_0)=p(s_1)$. The value $b(t)\in A_Q$ is a
single element. If $b(t)=a_Q(s_0)$ then, since $a_Q(s_1)\neq a_Q(s_0)$, we get
$b(p(s_1))=b(t)=a_Q(s_0)\neq a_Q(s_1)$, so $b$ errs at $s_1$. If instead $b(t)\neq a_Q(s_0)$
then $b(p(s_0))\neq a_Q(s_0)$ and $b$ errs at $s_0$. In either case $b$ fails at one of the two
states. Since $b$ was arbitrary, no map $\Rdt\to A_Q$ satisfies $a_Q=b\circ p$, which is the
negation of \cref{def:solvable}; this also follows from
\cref{thm:exact-solvability}\,\eqref{it:imp}, whose hypothesis is contradicted by
\eqref{eq:obstruction-hyp}.
\end{proof}

The structure of this argument deserves emphasis, because it is the pattern on which all the
necessity results of this paper are built, and because it is easy to counterfeit. We did not
assume that no readout-only construction exists and then declare completion necessary. We assumed
only \eqref{eq:obstruction-hyp} --- a fact about the pair $(p,a_Q)$, checkable in examples, and
verified concretely in \cref{sec:spectrum} --- and then \emph{derived} the failure of every
candidate $b$ by a counting argument: $b(t)$ is one element of $A_Q$, while correctness at $s_0$
and $s_1$ demands two different elements. Impossibility follows from variation of the
correct-answer map along a fiber, not from a stipulation of impossibility.

\begin{proposition}[Every readout-only construction is fiber-constant]\label{prop:fiberconst-form}
Let $X$ be any set and let $g\colon\Src\to X$ be a map that \emph{factors through the readout},
i.e.\ $g=h\circ p$ for some $h\colon\Rdt\to X$. Then $g$ is constant on every fiber of $p$.
Conversely, every map $\Src\to X$ constant on all fibers of $p$ factors through $p$, uniquely.
\end{proposition}

\begin{proof}
If $g=h\circ p$ and $s,s'\in F_t$, then $g(s)=h(t)=g(s')$. The converse and the uniqueness are
\cref{thm:exact-solvability} applied with $A_Q$ replaced by $X$ and $a_Q$ by $g$.
\end{proof}

\begin{proposition}[Auxiliary readout-computed variables cannot repair the
obstruction]\label{prop:auxiliary-useless}
Assume \eqref{eq:obstruction-hyp}. Let $W$ be any set and let $w\colon\Src\to W$ be an auxiliary
quantity that is \emph{computed from the readout alone}, meaning $w=\omega\circ p$ for some
$\omega\colon\Rdt\to W$. Then for every map
\begin{equation*}
c\colon\Rdt\times W\longrightarrow A_Q,
\end{equation*}
the induced construction $s\mapsto c(p(s),w(s))$ fails to equal $a_Q$ at $s_0$ or at $s_1$.
\end{proposition}

\begin{proof}
Define $b\colon\Rdt\to A_Q$ by $b(t):=c(t,\omega(t))$. Then for every $s\in\Src$,
\begin{equation*}
c\bigl(p(s),w(s)\bigr)=c\bigl(p(s),\omega(p(s))\bigr)=b(p(s)),
\end{equation*}
so the construction $s\mapsto c(p(s),w(s))$ is exactly $b\circ p$ and hence, by
\cref{prop:fiberconst-form}, constant on fibers. By \cref{cor:obstruction} it errs at $s_0$ or at
$s_1$.
\end{proof}

\Cref{prop:auxiliary-useless} closes an escape route that is attempted surprisingly often. One
cannot evade a fiber-variation obstruction by manufacturing additional quantities out of the
readout --- derived fields, filtered or smoothed versions of $t$, spectral transforms of $t$,
statistics of $t$, or any other post-processing. Whatever is computed from $t$ is a function of
$t$; composing functions of $t$ leaves one with a function of $t$; and a function of $t$ is
constant on $F_t$ while $a_Q$ is not. Genuine repair must inject information that is
\emph{not} a function of the readout. That is what \cref{sec:extensions} formalizes.

%===============================================================
\section{Active and inactive information loss}\label{sec:active}

We now sharpen the picture. \Cref{cor:obstruction} tells us that fiber variation of $a_Q$
obstructs. But in practice one usually knows something more specific: one knows \emph{which}
source-level structure the readout has forgotten, and one wants to know whether \emph{that
particular} forgotten structure is what the target depends on. This is the distinction between
information that is merely lost and information whose loss has consequences.

\begin{definition}[Source relation and its target factorization]\label{def:relation}
Let $r\colon\Src\to V$ be a map, called a \emph{source relation}, thought of as recording some
physically meaningful attribute of the reduced source state that the readout does not itself
determine. We say that $Q$ \emph{factors through} $(p,r)$ if there exists a map
\begin{equation}\label{eq:PhiQ}
\Phi_Q\colon \Rdt\times V\longrightarrow A_Q
\qquad\text{such that}\qquad
a_Q(s)=\Phi_Q\bigl(p(s),r(s)\bigr)\quad\text{for all }s\in\Src.
\end{equation}
\end{definition}

Equation \eqref{eq:PhiQ} says that the readout value together with the relation value suffices to
determine the correct answer. It is a modeling hypothesis about how much of the source the target
actually consults; the interesting content lies in what one may then conclude.

\begin{definition}[Activity and inactivity on a fiber]\label{def:active}
Assume \eqref{eq:PhiQ}, and let $t\in\Rdt$. The relation $r$ is \emph{active for $Q$ on the fiber
$F_t$} if there exist $s_0,s_1\in F_t$ with
\begin{equation}\label{eq:active}
\Phi_Q\bigl(t,r(s_0)\bigr)\neq\Phi_Q\bigl(t,r(s_1)\bigr).
\end{equation}
It is \emph{inactive for $Q$ on $F_t$} if no such pair exists, i.e.\ if the map
$\Phi_Q(t,\cdot)$ is constant on the \emph{realized set}
\begin{equation}\label{eq:realized}
r(F_t)=\{\, r(s) : s\in F_t \,\}\subseteq V .
\end{equation}
We say $r$ is \emph{globally inactive for $Q$} if it is inactive on every fiber.
\end{definition}

Two features of \cref{def:active} carry the weight of the section, and both are easy to get
wrong.

First, inactivity is a statement about the \emph{realized target values}, not about $r$. The
relation $r$ may vary wildly across $F_t$ --- $r(F_t)$ may be an enormous subset of $V$ --- and
$r$ may still be inactive for $Q$ on $F_t$, provided $\Phi_Q(t,\cdot)$ happens to be constant on
that subset. The readout has then forgotten a great deal, and none of it matters for $Q$. This
is the normal situation in effective descriptions and is not a pathology.

Second, activity is evaluated on the realized set $r(F_t)$, not on all of $V$. That
$\Phi_Q(t,\cdot)$ is nonconstant somewhere on $V$ is irrelevant if the source states in $F_t$
never realize the values at which it varies. The obstruction concerns configurations that
actually occur in the fiber.

\begin{theorem}[Active loss obstructs]\label{thm:active-obstructs}
Assume \eqref{eq:PhiQ}. If $r$ is active for $Q$ on some fiber $F_{t}$, then $Q$ is not exactly
solvable from the readout $p$.
\end{theorem}

\begin{proof}
Let $s_0,s_1\in F_{t}$ realize \eqref{eq:active}. Since $s_0,s_1\in F_{t}$ we have
$p(s_0)=t=p(s_1)$. By \eqref{eq:PhiQ},
\begin{equation*}
a_Q(s_0)=\Phi_Q\bigl(p(s_0),r(s_0)\bigr)=\Phi_Q\bigl(t,r(s_0)\bigr)
\neq
\Phi_Q\bigl(t,r(s_1)\bigr)=\Phi_Q\bigl(p(s_1),r(s_1)\bigr)=a_Q(s_1).
\end{equation*}
Thus $s_0,s_1$ satisfy the hypothesis \eqref{eq:obstruction-hyp} of \cref{cor:obstruction}, and
$Q$ is not exactly solvable from $p$. Indeed no map $b\colon\Rdt\to A_Q$ is correct at both.
\end{proof}

\begin{theorem}[Inactive loss does not obstruct]\label{thm:inactive-preserves}
Assume \eqref{eq:PhiQ}. If $r$ is globally inactive for $Q$ --- that is, if for every $t\in\Rdt$
the map $\Phi_Q(t,\cdot)$ is constant on the realized set $r(F_t)$ --- then $Q$ \emph{is} exactly
solvable from the readout $p$. The unique readout-only solution is
\begin{equation}\label{eq:inactive-solution}
\bar a_Q(t)=\Phi_Q\bigl(t,v_t\bigr)
\qquad\text{for any } v_t\in r(F_t).
\end{equation}
\end{theorem}

\begin{proof}
We verify condition \eqref{it:fibconst} of \cref{thm:exact-solvability}. Fix $t\in\Rdt$ and let
$s,s'\in F_t$. Then $r(s),r(s')\in r(F_t)$, and by global inactivity $\Phi_Q(t,\cdot)$ is
constant on $r(F_t)$; hence $\Phi_Q(t,r(s))=\Phi_Q(t,r(s'))$. Using \eqref{eq:PhiQ} and
$p(s)=p(s')=t$,
\begin{equation*}
a_Q(s)=\Phi_Q\bigl(t,r(s)\bigr)=\Phi_Q\bigl(t,r(s')\bigr)=a_Q(s').
\end{equation*}
So $a_Q$ is constant on $F_t$, for every $t$. By \cref{thm:exact-solvability}, $Q$ is exactly
solvable and the unique solution satisfies $\bar a_Q(t)=a_Q(s)$ for any $s\in F_t$; since
$F_t\neq\emptyset$ we may pick such an $s$ and set $v_t:=r(s)\in r(F_t)$, giving
\eqref{eq:inactive-solution}. The right-hand side of \eqref{eq:inactive-solution} is independent
of the choice of $v_t\in r(F_t)$, again by constancy of $\Phi_Q(t,\cdot)$ on $r(F_t)$.
\end{proof}

\begin{definition}[Structural versus consequential loss]\label{def:loss-types}
Let $r\colon\Src\to V$ be a source relation that does not factor through $p$; that is, $r$ is not
constant on some fiber. We call this \emph{structural information loss}: the readout genuinely
fails to determine $r$. If in addition $r$ is active for a target $Q$ on some fiber, we call the
loss \emph{consequential for $Q$}, or \emph{target-active}.
\end{definition}

\Cref{thm:active-obstructs,thm:inactive-preserves} should be read with their standing hypothesis
kept in view. Both presuppose that the nominated variable $r$ \emph{mediates} the target, in the
precise sense of \eqref{eq:PhiQ}: $a_Q(s)=\Phi_Q(p(s),r(s))$. Relative to such an $r$, the two
theorems say that activity of $r$ is exactly the condition under which the loss of $r$ obstructs
$Q$ --- active loss obstructs, and globally inactive loss does not. This is a statement about the
chosen mediating variable, not about information loss in the abstract, and it is only as sharp as
the choice of $r$.

The fully general criterion, carrying no such hypothesis, is the one already established in
\cref{thm:exact-solvability,cor:obstruction}: $Q$ is obstructed from the readout precisely when
$a_Q$ is nonconstant on some fiber of $p$. Every statement about activity is a refinement of that
criterion, obtained by localizing the fiber variation onto a particular mediating variable, and it
adds diagnostic content --- \emph{which} forgotten structure is responsible --- rather than logical
strength.

With that scoping in place, the moral stands. Structural loss without activity is invisible to
$Q$; it is what makes coarse descriptions viable rather than defective. It is therefore \emph{not}
a theorem of this framework --- and we do not assert --- that every noninjective readout requires
completion, nor that every lost source distinction must be restored. Whether a given loss must be
repaired is a question one can only answer after nominating a target, and the answer may
legitimately be no. The examples of \cref{sec:flat,sec:spectrum,sec:minimal-spectral} include both
verdicts for the same readout.

\begin{remark}\label{rem:activity-vs-descent}
It is worth recording the logical relation between activity of $r$ and descent of $r$. If $r$
itself descends along $p$ (is fiber-constant), then $r(F_t)$ is a singleton for each $t$ and $r$
is automatically inactive for every target that factors through $(p,r)$; the loss is not merely
inconsequential but nonexistent. The converse fails: $r$ can vary on fibers --- structural loss
--- while being inactive for $Q$. So activity implies structural loss, but structural loss does
not imply activity, and it is precisely this asymmetry that forbids the inference from
``noninjective'' to ``must be completed.''
\end{remark}

%===============================================================
\section{Minimal sound uncertainty without completion}\label{sec:uncertainty}

Suppose the obstruction is real: $a_Q$ genuinely varies on some fiber, so no readout-only
construction returns the correct answer everywhere. It does not follow that the readout is
useless for $Q$. It follows that the readout cannot deliver a \emph{single} answer, and the
honest question becomes: what is the most that the readout does license? The answer is a
set-valued rule, and there is a canonical least one.

\begin{definition}[Compatible-answer set]\label{def:compat}
For $t\in\Rdt$, the \emph{compatible-answer set} at $t$ is the image of the fiber under the
correct-answer map,
\begin{equation}\label{eq:compat}
\CompatSet_Q(t):=a_Q(F_t)=\{\, a_Q(s) : s\in\Src,\ p(s)=t \,\}\subseteq A_Q .
\end{equation}
Since $p$ is surjective, $F_t\neq\emptyset$ and hence $\CompatSet_Q(t)\neq\emptyset$.
\end{definition}

\begin{definition}[Set-valued rules and uniform soundness]\label{def:sound}
Let $\Pow(A_Q)$ denote the \emph{power set} of $A_Q$, i.e.\ the collection of all subsets of
$A_Q$. A \emph{set-valued rule} for $Q$ is a map
\begin{equation}\label{eq:rule}
C\colon\Rdt\longrightarrow\Pow(A_Q).
\end{equation}
The rule $C$ is \emph{uniformly sound} if it never excludes the truth:
\begin{equation}\label{eq:soundness}
a_Q(s)\in C\bigl(p(s)\bigr)\qquad\text{for every } s\in\Src .
\end{equation}
\end{definition}

\begin{remark}[$\Pow$ is the power set, not a space of probability measures]\label{rem:powerset}
We use the deliberately distinct symbol $\Pow(A_Q)$ to forestall a conflation that would destroy
the meaning of the results in this section. A set-valued rule assigns to each readout value a
\emph{subset} of admissible answers. It does not assign a probability distribution, a measure, a
density, or a credence. Nothing in \cref{def:sound} presupposes a $\sigma$-algebra on $A_Q$, a
measure on $\Src$, or a prior on the fibers, and no such structure is introduced anywhere in this
paper. See \cref{rem:not-probability} below for why this distinction is substantive and not
merely notational.
\end{remark}

\begin{theorem}[The compatible-answer set is the least uniformly sound
rule]\label{thm:minimal-sound}
Define $C_{\min}\colon\Rdt\to\Pow(A_Q)$ by $C_{\min}(t):=\CompatSet_Q(t)=a_Q(F_t)$. Then:
\begin{enumerate}
\item\label{it:cmin-sound} $C_{\min}$ is uniformly sound;
\item\label{it:cmin-least} for every uniformly sound rule $C\colon\Rdt\to\Pow(A_Q)$ and every
$t\in\Rdt$,
\begin{equation}\label{eq:cmin-least}
C_{\min}(t)\subseteq C(t).
\end{equation}
\end{enumerate}
Consequently $C_{\min}$ is the unique pointwise-least uniformly sound set-valued rule for $Q$
given the readout $p$.
\end{theorem}

\begin{proof}
\eqref{it:cmin-sound} Let $s\in\Src$ and put $t:=p(s)$. Then $s\in F_t$, so
$a_Q(s)\in a_Q(F_t)=C_{\min}(t)=C_{\min}(p(s))$. This is exactly \eqref{eq:soundness}.

\eqref{it:cmin-least} Let $C$ be uniformly sound and fix $t\in\Rdt$. Let
$\alpha\in C_{\min}(t)=a_Q(F_t)$ be arbitrary. By definition of the image there exists $s\in F_t$
with $a_Q(s)=\alpha$. Since $s\in F_t$ we have $p(s)=t$, and uniform soundness \eqref{eq:soundness}
applied to this $s$ gives
\begin{equation*}
\alpha=a_Q(s)\in C\bigl(p(s)\bigr)=C(t).
\end{equation*}
As $\alpha\in C_{\min}(t)$ was arbitrary, $C_{\min}(t)\subseteq C(t)$.

Uniqueness of the least element: if $C$ and $C'$ are both uniformly sound and both pointwise
least, then $C(t)\subseteq C'(t)$ and $C'(t)\subseteq C(t)$ for every $t$, so $C=C'$. Since
$C_{\min}$ is uniformly sound by \eqref{it:cmin-sound} and pointwise least by
\eqref{it:cmin-least}, it is the unique such rule.
\end{proof}

\begin{corollary}[Exactness as singleton-valuedness]\label{cor:singleton}
$Q$ is exactly solvable from the readout $p$ if and only if $\CompatSet_Q(t)$ is a singleton for
every $t\in\Rdt$. In that case the unique readout-only solution $\bar a_Q$ is characterized by
$\{\bar a_Q(t)\}=\CompatSet_Q(t)$.
\end{corollary}

\begin{proof}
$\CompatSet_Q(t)=a_Q(F_t)$ is a singleton for every $t$ iff $a_Q$ is constant on every fiber
$F_t$ (nonemptiness of $F_t$ gives that the image is nonempty, so ``singleton'' is equivalent to
``at most one element''), which by \cref{thm:exact-solvability} is equivalent to exact
solvability. If exact solvability holds, then for $s\in F_t$ we have
$a_Q(F_t)=\{a_Q(s)\}=\{\bar a_Q(p(s))\}=\{\bar a_Q(t)\}$.
\end{proof}

\begin{remark}[Set-valued uncertainty is not probability, and probability is not
recovery]\label{rem:not-probability}
$\CompatSet_Q(t)$ is the exact residual uncertainty that the readout partition induces about the
answer to $Q$ --- and it exists prior to, and independently of, any probabilistic structure. Two
consequences must be kept apart.

First, imposing a probability law on $\Src$, or on the fibers, changes the object under
discussion but does not shrink $\CompatSet_Q(t)$. A distribution on $A_Q$ can \emph{rank} the
compatible answers, declaring some more likely than others. It cannot make two distinct correct
answers into one answer, and it cannot identify which source state was in fact realized. We are
careful about how to phrase the constraint this places on the distribution. We do not claim that
an arbitrary probability law must assign positive mass to every element of $\CompatSet_Q(t)$; a
law is free to assign mass zero wherever it likes, and no full-support condition has been imposed
here. What is true is conditional: \emph{a probabilistic representation intended not to exclude
any source-compatible outcome must have support containing $\CompatSet_Q(t)$}. Under that
stipulation --- and it is a stipulation, not a theorem --- the support still contains both
$a_Q(s_0)$ and $a_Q(s_1)$ whenever $a_Q(s_0)\neq a_Q(s_1)$ and $p(s_0)=p(s_1)$, so the rule has
not told us which of $s_0,s_1$ obtained, and \cref{cor:obstruction} still forbids a correct
deterministic answer.

Second, and consequently, a probabilistic prediction answers a \emph{different target}. Replacing
$Q$ by the target ``what is the distribution of answers, given $t$ and given a specified law on
the fibers?'' is a legitimate move, but it is a change of question, and it is one that requires
additional input --- namely the law on the fibers, which is a piece of source-level information
not contained in $p$. It is not a solution to the original $Q$, and it should not be reported as
one. We return to this point in the dynamical setting in \cref{rem:stochastic}.
\end{remark}

%===============================================================
\section{Readout extensions and target sufficiency}\label{sec:extensions}

If the readout does not determine the answer, one may enlarge the readout. The following
definition says what an enlargement is, and the crucial point is that it need not be an
enlargement all the way back to the source.

\begin{definition}[Readout extension]\label{def:extension}
A \emph{readout extension} (or \emph{enrichment}) of $p\colon\Src\to\Rdt$ is a surjective map
\begin{equation}\label{eq:extension}
e\colon\Src\longrightarrow E
\qquad\text{together with a map}\qquad
u\colon E\longrightarrow\Rdt
\qquad\text{such that}\qquad
p=u\circ e .
\end{equation}
\end{definition}

\begin{lemma}[Basic properties of extensions]\label{lem:extension-basic}
Let $(e,E,u)$ be a readout extension of $p$. Then:
\begin{enumerate}
\item\label{it:u-unique} $u$ is uniquely determined by $e$ and $p$;
\item\label{it:u-surj} $u$ is surjective;
\item\label{it:refine} the partition of $\Src$ into fibers of $e$ refines the partition into
fibers of $p$; equivalently $e(s)=e(s')\implies p(s)=p(s')$.
\end{enumerate}
\end{lemma}

\begin{proof}
\eqref{it:u-unique} If $u,u'$ both satisfy $p=u\circ e=u'\circ e$, let $x\in E$. By surjectivity
of $e$ choose $s$ with $e(s)=x$; then $u(x)=u(e(s))=p(s)=u'(e(s))=u'(x)$.

\eqref{it:u-surj} Let $t\in\Rdt$. By surjectivity of $p$ choose $s$ with $p(s)=t$; then
$u(e(s))=p(s)=t$, so $t\in\Img(u)$.

\eqref{it:refine} If $e(s)=e(s')$ then $p(s)=u(e(s))=u(e(s'))=p(s')$.
\end{proof}

\Cref{lem:extension-basic}\,\eqref{it:refine} is the content of the word ``extension'': $E$ is at
least as informative as $\Rdt$, because anything determined by $E$ determines the $\Rdt$-value.
Equivalently, the equivalence relation $\sim_e$ induced by $e$ is contained in $\sim_p$. But $E$
need not be $\Src$, and $\sim_e$ need not be equality: $E$ sits somewhere between the readout and
the source, an intermediate quotient of $\Src$, or, in partition language, a refinement of the
readout partition that need not be the discrete partition. The two extreme cases are $E=\Rdt$
(with $e=p$, $u=\id$), which adds nothing, and $E=\Src$ (with $e=\id$, $u=p$), which retains
everything. The interesting extensions are strictly between, and the point of the next two
sections is that for a given target there is a canonical one, generally far closer to $\Rdt$ than
to $\Src$.

\begin{definition}[Sufficiency for a target]\label{def:sufficient}
A readout extension $e\colon\Src\to E$ is \emph{sufficient for $Q$}, or \emph{$Q$-sufficient}, if
there exists a map
\begin{equation}\label{eq:sufficiency}
b\colon E\longrightarrow A_Q
\qquad\text{such that}\qquad
a_Q=b\circ e .
\end{equation}
We also call a $Q$-sufficient readout extension a \emph{completion of $p$ for $Q$}.
\end{definition}

\begin{remark}[What ``completion'' means here]\label{rem:completion-meaning}
The word ``completion'' is used in \cref{def:sufficient} in a sense that must be separated from
several established and entirely different uses. It is \emph{not} metric completion, not
topological or uniform completion, not Cauchy completion, not the completion of a local ring, and
not the adjunction of arbitrary hidden variables posited to restore determinism. A completion in
the present sense is an \emph{enrichment of the readout that suffices for a nominated target}: a
surjection $e$ through which both $p$ and $a_Q$ factor. It is target-relative by construction; an
extension may complete $p$ for one target while remaining insufficient for another. And it is
constrained: \cref{thm:separation} shows that a completion cannot be arbitrary, and
\cref{thm:universal} shows that there is a least one.
\end{remark}

\begin{lemma}[Sufficiency criterion]\label{lem:sufficiency-criterion}
A readout extension $e\colon\Src\to E$ is $Q$-sufficient if and only if $a_Q$ is constant on
every fiber of $e$; and then $b$ in \eqref{eq:sufficiency} is unique.
\end{lemma}

\begin{proof}
Apply \cref{thm:exact-solvability} with $p$ replaced by the surjection $e$ and $\Rdt$ by $E$.
\end{proof}

Note the immediate sanity check: $e=\id_\Src$ (with $E=\Src$, $u=p$, $b=a_Q$) is always
$Q$-sufficient, so completions exist trivially. The theorems that follow are therefore not
existence theorems. They are theorems about \emph{how little} a completion can get away with, and
\emph{what it must contain no matter what}.

%===============================================================
\section{What every sufficient extension must retain}\label{sec:separation}

\begin{theorem}[Necessary separation]\label{thm:separation}
Let $e\colon\Src\to E$ be any $Q$-sufficient readout extension. Then $e$ separates every pair of
source states with different correct answers:
\begin{equation}\label{eq:separation}
a_Q(s_0)\neq a_Q(s_1)
\qquad\Longrightarrow\qquad
e(s_0)\neq e(s_1),
\qquad s_0,s_1\in\Src .
\end{equation}
\end{theorem}

\begin{proof}
We prove the contrapositive. Suppose $e(s_0)=e(s_1)=:x$. By $Q$-sufficiency there is
$b\colon E\to A_Q$ with $a_Q=b\circ e$. Then
\begin{equation*}
a_Q(s_0)=b\bigl(e(s_0)\bigr)=b(x)=b\bigl(e(s_1)\bigr)=a_Q(s_1),
\end{equation*}
since $b(x)$ is a single element of $A_Q$. Hence $e(s_0)=e(s_1)\implies a_Q(s_0)=a_Q(s_1)$, which
is \eqref{eq:separation}.
\end{proof}

\begin{corollary}[The critical case: fiber-mates with different answers]\label{cor:critical-pair}
Suppose $s_0,s_1\in\Src$ satisfy $p(s_0)=p(s_1)$ and $a_Q(s_0)\neq a_Q(s_1)$ --- that is, they
form an obstructing pair in the sense of \eqref{eq:obstruction-hyp}. Then every $Q$-sufficient
readout extension $e$ satisfies $e(s_0)\neq e(s_1)$, while $u(e(s_0))=u(e(s_1))$. In particular
$e$ is strictly finer than $p$: the extension is not merely allowed to distinguish $s_0$ from
$s_1$, it is \emph{compelled} to, and it must do so while remaining compatible with the readout
value they share.
\end{corollary}

\begin{proof}
The separation $e(s_0)\neq e(s_1)$ is \cref{thm:separation}. The identity
$u(e(s_i))=p(s_i)$ from \eqref{eq:extension}, together with $p(s_0)=p(s_1)$, gives
$u(e(s_0))=u(e(s_1))$. Since $e(s_0)\neq e(s_1)$ but $p$ does not separate $s_0,s_1$, the
partition $\sim_e$ is a strict refinement of $\sim_p$.
\end{proof}

It is useful to restate \cref{thm:separation} in partition language, where its content is
transparent. Let $\sim_e$, $\sim_p$, and $\sim_{a_Q}$ denote the equivalence relations on $\Src$
induced respectively by $e$, by $p$, and by $a_Q$ (so $s\sim_{a_Q}s'$ iff $a_Q(s)=a_Q(s')$).
\Cref{lem:extension-basic}\,\eqref{it:refine} says $\sim_e\;\subseteq\;\sim_p$, and
\cref{thm:separation} says $\sim_e\;\subseteq\;\sim_{a_Q}$. Thus:

\begin{quote}
\emph{A readout extension is sufficient for $Q$ precisely when its induced partition refines both
the readout partition and the answer partition.}
\end{quote}

Indeed the converse of \cref{thm:separation} holds as well and is worth recording, since together
they characterize sufficiency exactly.

\begin{proposition}[Separation characterizes sufficiency]\label{prop:separation-iff}
A readout extension $e\colon\Src\to E$ is $Q$-sufficient if and only if \eqref{eq:separation}
holds.
\end{proposition}

\begin{proof}
Necessity is \cref{thm:separation}. For sufficiency, assume \eqref{eq:separation}. Its
contrapositive reads $e(s)=e(s')\implies a_Q(s)=a_Q(s')$, i.e.\ $a_Q$ is constant on every fiber
of $e$. By \cref{lem:sufficiency-criterion}, $e$ is $Q$-sufficient.
\end{proof}

We emphasize that \cref{thm:separation} does \emph{not} say that $e$ must separate every pair of
distinct source states, i.e.\ that $e$ must be injective, i.e.\ that the source must be
reconstructed. It says only that $e$ must separate pairs \emph{whose answers differ}. Pairs with
equal answers may be freely collapsed. That is the loophole that \cref{sec:joint} exploits to
construct a completion that is often dramatically smaller than $\Src$, and it is why a
target-relative theory of completion is not a disguised demand for full reconstruction.

%===============================================================
\section{The joint target image and universal minimality}\label{sec:joint}

\Cref{thm:separation} gives a lower bound on the distinctions a completion must carry. We now
show that the bound is attained within the category of readout extensions defined above: there
is a completion carrying exactly those distinctions and no others, and it receives a canonical
map from every other completion. Attainment here is a statement about the abstract construction
$E_Q=\Img(p,a_Q)$; whether a given physical situation supplies an operational realization of that
completion is a separate matter, taken up after \cref{thm:universal}.

\begin{definition}[Joint target image]\label{def:joint}
Define
\begin{equation}\label{eq:jQ}
j_Q\colon\Src\longrightarrow\Rdt\times A_Q,
\qquad
j_Q(s):=\bigl(p(s),a_Q(s)\bigr),
\end{equation}
and let
\begin{equation}\label{eq:EQ}
E_Q:=\Img(j_Q)=\bigl\{\,(p(s),a_Q(s)) : s\in\Src \,\bigr\}\subseteq\Rdt\times A_Q .
\end{equation}
We regard $j_Q$ as a map onto its image, $j_Q\colon\Src\to E_Q$, and write
$\pi_\Rdt\colon E_Q\to\Rdt$ and $\pi_A\colon E_Q\to A_Q$ for the restrictions to $E_Q$ of the two
coordinate projections of $\Rdt\times A_Q$.
\end{definition}

\begin{proposition}[$E_Q$ is a $Q$-sufficient readout extension]\label{prop:EQ-sufficient}
The map $j_Q\colon\Src\to E_Q$ is a surjective readout extension of $p$, with structure map
$u=\pi_\Rdt$, and it is sufficient for $Q$, with $b=\pi_A$. Explicitly,
\begin{equation}\label{eq:EQ-factorizations}
p=\pi_\Rdt\circ j_Q,
\qquad\qquad
a_Q=\pi_A\circ j_Q .
\end{equation}
Moreover the fibers of $j_Q$ are exactly the sets $F_t\cap a_Q^{-1}(\alpha)$ for
$(t,\alpha)\in E_Q$.
\end{proposition}

\begin{proof}
$j_Q\colon\Src\to E_Q$ is surjective by construction, since $E_Q$ was defined as the image of
$j_Q$. No nonemptiness hypothesis is needed: if $\Src=\emptyset$ then $E_Q=\emptyset$ as well, and
the identities below hold vacuously.

For \eqref{eq:EQ-factorizations}, let $s\in\Src$. Then $j_Q(s)=(p(s),a_Q(s))\in E_Q$, and
applying the coordinate projections,
\begin{equation*}
\pi_\Rdt\bigl(j_Q(s)\bigr)=\pi_\Rdt\bigl(p(s),a_Q(s)\bigr)=p(s),
\qquad
\pi_A\bigl(j_Q(s)\bigr)=\pi_A\bigl(p(s),a_Q(s)\bigr)=a_Q(s).
\end{equation*}
The first identity says $p=\pi_\Rdt\circ j_Q$, so $(j_Q,E_Q,\pi_\Rdt)$ is a readout extension in
the sense of \cref{def:extension}. The second says $a_Q=\pi_A\circ j_Q$, so $j_Q$ is
$Q$-sufficient in the sense of \cref{def:sufficient}, with $b=\pi_A$.

For the fibers: $j_Q(s)=(t,\alpha)$ iff $p(s)=t$ and $a_Q(s)=\alpha$, i.e.\ iff
$s\in F_t\cap a_Q^{-1}(\alpha)$.
\end{proof}

We have thus verified directly, and not by appeal to any universal property, that $E_Q$ does the
job: it determines the readout value and it determines the correct answer. The next theorem says
it does nothing more.

\begin{theorem}[Universal factorization: $E_Q$ is the coarsest $Q$-sufficient
extension]\label{thm:universal}
Let $e\colon\Src\to E$ be \emph{any} surjective readout extension of $p$ that is sufficient for
$Q$, with structure map $u\colon E\to\Rdt$ satisfying $p=u\circ e$ and solution map
$b\colon E\to A_Q$ satisfying $a_Q=b\circ e$. Then there is a \emph{unique} map
\begin{equation}\label{eq:phiE}
\varphi_E\colon E\longrightarrow E_Q
\qquad\text{such that}\qquad
j_Q=\varphi_E\circ e,
\end{equation}
namely
\begin{equation}\label{eq:phiE-formula}
\varphi_E(x)=\bigl(u(x),\,b(x)\bigr),
\qquad x\in E,
\end{equation}
and $\varphi_E$ is surjective. Moreover $\varphi_E$ intertwines the structure maps and the
solution maps:
\begin{equation}\label{eq:phiE-intertwine}
\pi_\Rdt\circ\varphi_E=u,
\qquad\qquad
\pi_A\circ\varphi_E=b .
\end{equation}
\end{theorem}

\begin{proof}
We proceed in five steps.

\emph{Step 1: $u$ and $b$ are uniquely determined by $e$.} This is
\cref{lem:extension-basic}\,\eqref{it:u-unique} for $u$, and \cref{lem:sufficiency-criterion} for
$b$; in both cases uniqueness follows from surjectivity of $e$. We record the argument for $b$
once more, since it is used below: if $b,b'$ both satisfy $a_Q=b\circ e=b'\circ e$ and $x\in E$,
choose $s$ with $e(s)=x$; then $b(x)=b(e(s))=a_Q(s)=b'(e(s))=b'(x)$. Hence the formula
\eqref{eq:phiE-formula} is unambiguous.

\emph{Step 2: $\varphi_E$ is well defined as a map into $E_Q$.} Define
$\psi\colon E\to\Rdt\times A_Q$ by $\psi(x):=(u(x),b(x))$; this is manifestly a well-defined map
into the product. We must check that its values lie in the subset $E_Q$. Let $x\in E$. By
surjectivity of $e$ there is $s\in\Src$ with $e(s)=x$. Then
\begin{equation}\label{eq:psi-value}
\psi(x)=\bigl(u(e(s)),\,b(e(s))\bigr)=\bigl(p(s),\,a_Q(s)\bigr)=j_Q(s)\in E_Q,
\end{equation}
using $p=u\circ e$ and $a_Q=b\circ e$. Hence $\psi$ takes values in $E_Q$, and we may define
$\varphi_E\colon E\to E_Q$ by $\varphi_E:=\psi$, i.e.\ by \eqref{eq:phiE-formula}. (Note that
\eqref{eq:psi-value} did not require the choice of $s$ to be canonical: $\psi(x)$ was already
defined without reference to $s$, and $s$ was used only to certify membership in $E_Q$. So no
well-definedness issue arises from the choice of preimage.)

\emph{Step 3: the factorization $j_Q=\varphi_E\circ e$.} Let $s\in\Src$. Then by
\eqref{eq:phiE-formula} and the two factorizations through $e$,
\begin{equation*}
\varphi_E\bigl(e(s)\bigr)
=\bigl(u(e(s)),\,b(e(s))\bigr)
=\bigl(p(s),\,a_Q(s)\bigr)
=j_Q(s).
\end{equation*}
Since $s$ was arbitrary, $j_Q=\varphi_E\circ e$.

\emph{Step 4: $\varphi_E$ is surjective.} Let $y\in E_Q$. By \eqref{eq:EQ} there is $s\in\Src$
with $j_Q(s)=y$. Put $x:=e(s)\in E$. By Step 3, $\varphi_E(x)=\varphi_E(e(s))=j_Q(s)=y$. Hence
$\varphi_E$ is onto $E_Q$.

\emph{Step 5: uniqueness of $\varphi_E$.} Suppose $\varphi,\varphi'\colon E\to E_Q$ both satisfy
$j_Q=\varphi\circ e=\varphi'\circ e$. Let $x\in E$. By surjectivity of $e$ choose $s\in\Src$ with
$e(s)=x$. Then
\begin{equation*}
\varphi(x)=\varphi\bigl(e(s)\bigr)=j_Q(s)=\varphi'\bigl(e(s)\bigr)=\varphi'(x).
\end{equation*}
As $x$ was arbitrary, $\varphi=\varphi'$. This is precisely the point at which surjectivity of
$e$ --- part of \cref{def:extension} --- is essential: without it, $\varphi_E$ would be
unconstrained on $E\setminus\Img(e)$.

Finally, \eqref{eq:phiE-intertwine} is immediate from \eqref{eq:phiE-formula}: for $x\in E$,
$\pi_\Rdt(\varphi_E(x))=\pi_\Rdt(u(x),b(x))=u(x)$ and
$\pi_A(\varphi_E(x))=\pi_A(u(x),b(x))=b(x)$.
\end{proof}

\begin{corollary}[Information ordering]\label{cor:ordering}
Every $Q$-sufficient readout extension $e\colon\Src\to E$ carries at least the distinctions
encoded by the pair $(p,a_Q)$: for all $s_0,s_1\in\Src$,
\begin{equation}\label{eq:ordering}
j_Q(s_0)\neq j_Q(s_1)
\qquad\Longrightarrow\qquad
e(s_0)\neq e(s_1).
\end{equation}
Equivalently, $\sim_e\;\subseteq\;\sim_{j_Q}$: the partition induced by $e$ refines the partition
induced by $j_Q$. Thus $E_Q$ is the \emph{coarsest} readout extension sufficient for $Q$.
\end{corollary}

\begin{proof}
By \cref{thm:universal}, $j_Q=\varphi_E\circ e$. If $e(s_0)=e(s_1)$ then
$j_Q(s_0)=\varphi_E(e(s_0))=\varphi_E(e(s_1))=j_Q(s_1)$, which is the contrapositive of
\eqref{eq:ordering}. Hence $\sim_e\subseteq\sim_{j_Q}$. Since $j_Q$ is itself $Q$-sufficient
(\cref{prop:EQ-sufficient}), it is a $Q$-sufficient extension whose partition is refined by that
of every other, i.e.\ it is coarsest among $Q$-sufficient extensions.
\end{proof}

Several points about \cref{thm:universal} and \cref{cor:ordering} require care, and we state them
explicitly.

\emph{The direction of the universal map.} The canonical map runs
\begin{equation*}
E\longrightarrow E_Q,
\end{equation*}
from an arbitrary sufficient extension \emph{to} the joint target image --- not from $E_Q$ to
$E$. This direction is forced and is not a convention. A surjection $E\to E_Q$ expresses that $E$
is at least as fine as $E_Q$, since it says the $E$-value determines the $E_Q$-value. There is in
general \emph{no} map $E_Q\to E$ over $\Src$: $E$ may retain distinctions that $E_Q$ has
discarded --- distinctions irrelevant to $Q$ but possibly relevant to other targets --- and no
map out of $E_Q$ could recover them. Reversing the arrow would assert the opposite of what is
true, namely that $E_Q$ is the finest rather than the coarsest sufficient extension.

\emph{$E_Q$ does not reconstruct the source.} The joint target image retains exactly two things:
the readout value $p(s)$, and the correct answer $a_Q(s)$. It discards everything else, and in
particular it identifies any two source states that agree on both. This is usually a massive
collapse. In \cref{sec:minimal-spectral} the target is the spectrum of a charged Laplacian, and
$E_Q$ turns out to be a one-dimensional interval's worth of data --- the inversion quotient of
holonomy --- even though the source moduli space $\Uone$ is strictly finer and the raw
connections finer still. Nothing in \cref{thm:universal} says or implies that the source is
recoverable; on the contrary, \cref{cor:ordering} says the source is generally \emph{not}
recoverable from $E_Q$, because $\id_\Src$ is one of the sufficient extensions $e$ and the
canonical map goes $\Src\to E_Q$, not back.

\emph{One target at a time.} \Cref{thm:universal} is a statement about a single correct-answer map
$a_Q$. If one nominates a family of targets, the coarsest extension sufficient for all of them
simultaneously must, by \cref{thm:separation} applied to each member, separate every pair of source
states on which \emph{some} member's answer differs. It is therefore at least as fine as the
completion required by each individual member, and it may be strictly finer. It need not be:
should the remaining members contribute no distinctions beyond those already forced by one of
them, the family completion coincides with that member's individual completion. Which case obtains
depends on the family, and the construction in general requires machinery we do not develop here;
that is the subject of \cite{ShadowIII}. Within the present paper, ``minimal'' always
means ``minimal for the nominated $Q$,'' and we will see in \cref{sec:minimal-spectral} that the
minimal completion for the spectral target is strictly insufficient for a transport target on the
same source space.

\emph{What the universal property does and does not supply.} \Cref{thm:universal} establishes
that $E_Q$ is the coarsest sufficient extension in the category of set-theoretic source
quotients and readout extensions defined in this paper: among all $Q$-sufficient extensions, it
is the one carrying the fewest distinctions, and every other maps canonically onto it. Because
the construction of $E_Q=\Img(p,a_Q)$ refers explicitly to the correct-answer map $a_Q$, the
theorem is an information-theoretic characterization of the least distinctions any sufficient
extension must carry; it identifies the floor abstractly, from the data $(p,a_Q)$. It does not,
however, furnish an operational procedure for extracting $a_Q$ from the original readout $p$ ---
indeed \cref{cor:obstruction,prop:auxiliary-useless} show that no such extraction from $p$ alone
exists when $Q$ is obstructed --- and it does not by itself exhibit a physically measurable,
experimentally accessible, or independently specified realization of $E_Q$. A concrete physical
repair is a separate task: it must independently provide an extension $e\colon\Src\to E$ and an
answer map $b\colon E\to A_Q$ through which $a_Q$ factors, and must justify that $e$ and $b$ are
available in the situation at hand. The universal property then locates any such repair relative
to the floor --- it certifies that $e$ carries at least the distinctions of $E_Q$ --- but the
provision and justification of $e$ and $b$ lie outside what the abstract quotient construction
delivers. The flat-connection example of \cref{sec:flat,sec:minimal-spectral} supplies exactly
such an independently specified extension, the charged spectrum, and it is that concrete
construction, not \cref{thm:universal} alone, that makes the repair physical.

%===============================================================
\section{Flat \texorpdfstring{$\Uone$}{U(1)} connections on the circle}\label{sec:flat}

We now instantiate the framework in a setting where every ingredient --- the reduction by gauge,
the deliberately restricted readout, the fiber, the target, and the minimal completion --- is
explicit and computable.

\subsection{The reduced moduli space}

Let $S^1$ be a circle of circumference $L>0$, with coordinate $x\in[0,L]$ and endpoints
identified, $x\sim x+L$. Consider the trivial principal $\Uone$-bundle over $S^1$ and its
connections, described in the standard way by real one-forms $A$ on $S^1$; the associated
covariant derivative on a unit-charge complex scalar is $\partial_x-iA_x$, and the curvature is
the two-form $F=dA$.

We state at once the feature that makes this example the sharpest possible illustration of the
theory, rather than treating it as an obstacle to be maneuvered around. \emph{Every two-form on a
one-dimensional manifold vanishes identically}, since $\Lambda^2 T^*_xS^1=0$ for dimensional
reasons. Hence $F=dA=0$ for \emph{every} connection $A$ on $S^1$, with no flatness hypothesis
required and no special choice of representative involved. The curvature readout on this space is
therefore identically zero as a matter of dimension.

This is not a degeneracy to be apologized for; it is exactly the point. Curvature is a
\emph{local} invariant, and on the circle it is vacuous. Holonomy is a \emph{global} invariant,
and on the circle it is not: connections with identically vanishing field strength are
nonetheless distinguished, gauge-invariantly, by the phase a charge accumulates on traversing the
loop. The circle is thus the cleanest available laboratory for a global-holonomy-versus-local-
curvature separation, and that separation is precisely what \cref{sec:spectrum,sec:minimal-spectral}
will show a physical target can detect.

Consider now the one-parameter family of connection one-forms
\begin{equation}\label{eq:Atheta}
A_\theta=\frac{\theta}{L}\,dx,
\qquad
\theta\in\Rr .
\end{equation}
Their field strength vanishes,
\begin{equation}\label{eq:Ftheta}
F_\theta=dA_\theta=d\!\left(\frac{\theta}{L}\right)\wedge dx=0,
\end{equation}
which holds for the two independent reasons just discussed: the coefficient $\theta/L$ is
constant, and, more fundamentally, any two-form on $S^1$ is zero regardless. Every connection on
the trivial bundle over $S^1$ is gauge-equivalent to exactly one such $A_\theta$ modulo the
identifications described below; this is the standard parametrization of the moduli space of
$\Uone$ connections on a circle \cite{Kobayashi,NakaharaBook,BaezMuniain}. Since all connections
here are flat, we retain the customary name \emph{moduli space of flat connections} for it.

\subsection{Gauge convention and the reduction}

Gauge transformations act by
\begin{equation}\label{eq:gauge}
A\longmapsto A+d\lambda ,
\end{equation}
where $g=e^{i\lambda}\colon S^1\to\Uone$ is a smooth, single-valued map. Single-valuedness of $g$,
not of $\lambda$, is the physical requirement, and it is exactly this that produces the
identification we need. Writing $\lambda(x)=2\pi k x/L$ with $k\in\Zz$, the map
\begin{equation}\label{eq:largegauge}
g_k(x)=\exp\!\left(\frac{2\pi i k x}{L}\right)
\end{equation}
is smooth and single-valued on $S^1$, since $g_k(x+L)=g_k(x)e^{2\pi i k}=g_k(x)$, even though
$\lambda$ itself increases by $2\pi k$ around the circle. Such $g_k$ are the \emph{large gauge
transformations}; they are not connected to the identity, and they act on \eqref{eq:Atheta} by
\begin{equation}\label{eq:largegauge-action}
A_\theta\longmapsto A_\theta+d\lambda=\frac{\theta}{L}dx+\frac{2\pi k}{L}dx=A_{\theta+2\pi k}.
\end{equation}
Consequently the parameters $\theta$ and $\theta+2\pi k$ describe \emph{the same physical
configuration}, and the honest parameter space of flat connections is the quotient
\begin{equation}\label{eq:theta-quotient}
\theta\in\Rr/2\pi\Zz .
\end{equation}
Small gauge transformations, with $\lambda$ single-valued, change $A_\theta$ by an exact form
$d\lambda$ with $\oint d\lambda=0$; these do not alter $\theta$ mod $2\pi$ and, after fixing the
constant gauge \eqref{eq:Atheta}, do not move us within the family at all.

We may now define the reduced source space of this example in accordance with the standing
convention of \cref{sec:setup}:
\begin{equation}\label{eq:Sflat}
\Src_{\mathrm{flat}}:=\{\text{flat }\Uone\text{ connections on }S^1\}\big/\{\text{gauge}\}
\;\cong\;\Rr/2\pi\Zz\;\cong\;\Uone .
\end{equation}
It is essential to the correctness of everything below that the identification
$\theta\sim\theta+2\pi k$ has \emph{already been performed}. Distinct points of
$\Src_{\mathrm{flat}}$ --- distinct values of $e^{i\theta}$ --- are gauge-inequivalent physical
configurations, not gauge copies of one another. Any obstruction we exhibit between two such
points is therefore an obstruction between genuinely different physical states, and cannot be
dissolved by a change of gauge.

\subsection{Holonomy}

The gauge-invariant content of $A_\theta$ is its holonomy around the circle,
\begin{equation}\label{eq:holonomy}
h(\theta)=\exp\!\left(i\oint_{S^1}A_\theta\right)
=\exp\!\left(i\int_0^L\frac{\theta}{L}\,dx\right)
=e^{i\theta}\in\Uone .
\end{equation}
Under a large gauge transformation \eqref{eq:largegauge-action}, $h\mapsto e^{i(\theta+2\pi k)}
=e^{i\theta}$: the holonomy is invariant, as it must be. Under a small gauge transformation the
integral $\oint d\lambda$ vanishes and $h$ is again unchanged. So $h$ is a well-defined function
on $\Src_{\mathrm{flat}}$, and by \eqref{eq:Sflat} it is in fact a bijection
$\Src_{\mathrm{flat}}\to\Uone$. We will accordingly use $e^{i\theta}$, or $\theta$ mod $2\pi$,
as the coordinate on the reduced source space.

That the holonomy is physically consequential --- that it is not a gauge artifact but an
observable feature of the configuration --- is the content of the Aharonov--Bohm effect
\cite{AharonovBohm,WuYang}, and it is worth stating the geometry of that effect accurately.

In the standard physical realization, the apparatus consists of a solenoid or flux tube together
with the region exterior to it. The charged particle is confined to the \emph{accessible region},
which excludes the flux-carrying interior and is therefore multiply connected; on that accessible
region, and only there, the connection is flat and the field strength vanishes. The excluded
region does carry field: the flux $\Phi$ threading the solenoid is nonzero, and it is exactly that
flux which fixes the phase. It would be wrong to say that the electromagnetic field vanishes
everywhere in the apparatus --- it does not --- and we do not say so. What is true, and what
matters, is that the interference phase measured by the particle is determined by the
gauge-invariant holonomy of the connection around a loop in the accessible region, a loop along
which the particle never encounters any field.

The circle $S^1$ of this section models such an accessible loop. Its states $e^{i\theta}$ are
distinguished by holonomy alone, with no local field anywhere along them, which is the physical
warrant for treating $\Src_{\mathrm{flat}}\cong\Uone$ as a space of genuinely distinct physical
states rather than as gauge copies of a single trivial one.

\subsection{The curvature readout}

Define the readout to be the field strength:
\begin{equation}\label{eq:pF}
p_F\colon\Src_{\mathrm{flat}}\longrightarrow\{0\},
\qquad
p_F(\theta)=F_\theta=0 .
\end{equation}
By \eqref{eq:Ftheta} this is well defined and, trivially, surjective onto the one-point set
$\{0\}$: the curvature vanishes on every connection, for the dimensional reason recorded above,
so $\{0\}$ is the full set of curvature values realized. Its single fiber is everything:
\begin{equation}\label{eq:pF-fiber}
F_{0}=p_F^{-1}(0)=\Src_{\mathrm{flat}}\cong\Uone .
\end{equation}
This is about as lossy as a readout can be. It is also, for exactly that reason, an ideal test
case: any target whose correct answer is nonconstant on $\Uone$ is obstructed by $p_F$, and any
target whose correct answer is constant on $\Uone$ is not.

\begin{remark}[The curvature readout is a deliberate restriction, not a claim about
physics]\label{rem:curvature-restricted}
We are not asserting that curvature exhausts what can be observed of a $\Uone$ gauge field. It
plainly does not: in the Aharonov--Bohm interference experiment the measured interference phase is
determined by the gauge-invariant holonomy, and an experiment sensitive to holonomy would supply
a strictly finer readout than $p_F$ --- indeed
the finest possible one on $\Src_{\mathrm{flat}}$, since $h$ is injective there. The map $p_F$ is
the readout \emph{under examination}. It is nominated, not discovered, and it is nominated
precisely because its extreme degeneracy makes the target-relative machinery of
\cref{sec:exact,sec:active,sec:separation,sec:joint} maximally visible. A reader who objects that
``real physics sees the holonomy'' is agreeing with the analysis: enlarging the readout to
include holonomy is exactly the move that \cref{sec:extensions} formalizes, and
\cref{sec:minimal-spectral} shows that for the spectral target one need not even go that far.
\end{remark}

\begin{example}[A target unobstructed by the curvature readout]\label{ex:unobstructed}
Take $A_Q=\{\text{yes},\text{no}\}$ and $a_Q(\theta)=\text{yes}$ iff the connection is flat. Then
$a_Q$ is constant on $\Src_{\mathrm{flat}}$, hence constant on the fiber $F_0$, hence by
\cref{thm:exact-solvability} exactly solvable from $p_F$: the readout-only solution is
$\bar a_Q(0)=\text{yes}$. The curvature readout, degenerate as it is, answers this target
perfectly. The next two sections exhibit a target on the same source space, with the same
readout, for which it fails.
\end{example}

%===============================================================
\section{The charged covariant Laplacian and the curvature obstruction}\label{sec:spectrum}

\subsection{Operator, domain, and spectrum}

Fix $\theta\in\Rr$ representing a point of $\Src_{\mathrm{flat}}$, and consider a complex scalar
field of unit charge on $S^1$ minimally coupled to $A_\theta$. The gauge-covariant derivative is
\begin{equation}\label{eq:Dx}
\Dx:=\partial_x-\frac{i\theta}{L},
\end{equation}
and the covariant Laplacian, or gauged kinetic operator, is
\begin{equation}\label{eq:Htheta}
H_\theta:=-\Dx^{\,2}
=-\left(\partial_x-\frac{i\theta}{L}\right)^{\!2}.
\end{equation}

We specify the operator precisely. Let
\begin{equation}\label{eq:Hilbert}
\mathcal{H}:=L^2(S^1)=L^2\bigl([0,L],dx\bigr)
\quad\text{with periodic identification } x\sim x+L,
\end{equation}
with inner product $\langle\psi,\phi\rangle=\int_0^L\overline{\psi(x)}\phi(x)\,dx$, and take the
operator domain to be the periodic Sobolev space
\begin{equation}\label{eq:domain}
\mathcal{D}(H_\theta):=H^2_{\mathrm{per}}(S^1)
=\Bigl\{\psi\in L^2(S^1):\ \sum_{n\in\Zz}(1+n^2)^2\,|\hat\psi(n)|^2<\infty\Bigr\},
\end{equation}
where $\hat\psi(n)$ are the Fourier coefficients of $\psi$ with respect to the orthonormal basis
introduced below. Equivalently, $\mathcal D(H_\theta)$ consists of $\psi\in L^2$ whose
distributional derivatives up to order two lie in $L^2$ and which, together with their first
derivatives, satisfy the periodic boundary conditions $\psi(0)=\psi(L)$, $\psi'(0)=\psi'(L)$;
we write $H^2_{\mathrm{per}}(S^1)$ for this space.
With this domain, $H_\theta$ is self-adjoint, and the cleanest justification is Fourier
diagonalization rather than a perturbative argument. Let $\mathcal F\colon L^2(S^1)\to
\ell^2(\Zz)$, $(\mathcal F\psi)(n)=\hat\psi(n)=\langle\psi_n,\psi\rangle$, be the unitary Fourier
transform associated with the orthonormal basis $\{\psi_n\}_{n\in\Zz}$ of \eqref{eq:psin} below.
A direct computation, carried out in \cref{prop:spectrum}, gives $H_\theta\psi_n=
\lambda_n(\theta)\psi_n$ with $\lambda_n(\theta)=\bigl((2\pi n-\theta)/L\bigr)^2$. Hence
$\mathcal F$ conjugates $H_\theta$ into the \emph{multiplication operator} by the real sequence
$\bigl(\lambda_n(\theta)\bigr)_{n\in\Zz}$ on $\ell^2(\Zz)$:
\begin{equation}\label{eq:diagonalized}
\bigl(\mathcal F H_\theta\psi\bigr)(n)=\lambda_n(\theta)\,\hat\psi(n),
\qquad
\mathcal F\bigl(\mathcal D(H_\theta)\bigr)
=\Bigl\{\,\hat\psi\in\ell^2(\Zz):\ \sum_{n\in\Zz}\lambda_n(\theta)^2\,|\hat\psi(n)|^2<\infty\,\Bigr\}.
\end{equation}
Multiplication by a real sequence, taken on its maximal domain, is self-adjoint; and since
$\lambda_n(\theta)\asymp (2\pi n/L)^2$ as $|n|\to\infty$, the maximal domain in
\eqref{eq:diagonalized} coincides with $\mathcal F\bigl(H^2_{\mathrm{per}}(S^1)\bigr)$, so
$\mathcal D(H_\theta)=H^2_{\mathrm{per}}(S^1)$ as claimed. Because the multiplier
$\lambda_n(\theta)\to+\infty$, the operator has compact resolvent and purely discrete spectrum
\cite{ReedSimonI,ReedSimonIV}. No relative-boundedness estimate is needed, and none is
claimed: the first-order term $\tfrac{2i\theta}{L}\partial_x$ is \emph{not} bounded on $L^2(S^1)$,
and the diagonalization above avoids any appeal to such a statement.

\subsection{The spectrum}

\begin{proposition}[Eigenvalues and multiplicities]\label{prop:spectrum}
The functions
\begin{equation}\label{eq:psin}
\psi_n(x)=\frac{1}{\sqrt{L}}\,e^{2\pi i n x/L},
\qquad n\in\Zz,
\end{equation}
form an orthonormal basis of $\mathcal{H}$ consisting of eigenvectors of $H_\theta$, and
\begin{equation}\label{eq:Dpsin}
\Dx\psi_n=\frac{i\,(2\pi n-\theta)}{L}\,\psi_n,
\qquad\qquad
H_\theta\psi_n=\lambda_n(\theta)\,\psi_n
\quad\text{with}\quad
\lambda_n(\theta)=\left(\frac{2\pi n-\theta}{L}\right)^{\!2}.
\end{equation}
Consequently the spectrum of $H_\theta$ is purely discrete and, \emph{counted with multiplicity},
is the multiset
\begin{equation}\label{eq:spec}
\Spec(H_\theta)=\Bigl\{\!\!\Bigl\{\,\lambda_n(\theta)=\Bigl(\tfrac{2\pi n-\theta}{L}\Bigr)^{2}
\ :\ n\in\Zz \,\Bigr\}\!\!\Bigr\},
\end{equation}
where $\{\!\{\cdot\}\!\}$ denotes a multiset: the eigenvalue $\lambda$ occurs in
$\Spec(H_\theta)$ with multiplicity
\begin{equation}\label{eq:multiplicity}
m_\theta(\lambda)=\#\bigl\{\, n\in\Zz : \lambda_n(\theta)=\lambda \,\bigr\}.
\end{equation}
\end{proposition}

\begin{proof}
Orthonormality and completeness of $\{\psi_n\}$ in $L^2(S^1)$ is the classical Fourier basis
theorem. Each $\psi_n$ is smooth and periodic, hence lies in $H^2(S^1)=\mathcal D(H_\theta)$.
Compute
\begin{equation*}
\Dx\psi_n=\partial_x\psi_n-\frac{i\theta}{L}\psi_n
=\frac{2\pi i n}{L}\psi_n-\frac{i\theta}{L}\psi_n
=\frac{i(2\pi n-\theta)}{L}\,\psi_n,
\end{equation*}
which is the first identity in \eqref{eq:Dpsin}. Applying $\Dx$ once more,
\begin{equation*}
\Dx^{\,2}\psi_n=\frac{i(2\pi n-\theta)}{L}\,\Dx\psi_n
=\left(\frac{i(2\pi n-\theta)}{L}\right)^{\!2}\psi_n
=-\left(\frac{2\pi n-\theta}{L}\right)^{\!2}\psi_n,
\end{equation*}
so $H_\theta\psi_n=-\Dx^2\psi_n=\bigl(\tfrac{2\pi n-\theta}{L}\bigr)^2\psi_n$, which is the
second identity. Since $\{\psi_n\}_{n\in\Zz}$ is a complete orthonormal system of eigenvectors with real
eigenvalues $\lambda_n(\theta)\to+\infty$ as $|n|\to\infty$, the Fourier transform $\mathcal F$
carries $H_\theta$ with domain \eqref{eq:domain} onto multiplication by
$(\lambda_n(\theta))_{n\in\Zz}$ on its maximal domain, as in \eqref{eq:diagonalized}. That
operator is self-adjoint with compact resolvent, hence so is $H_\theta$, and its spectrum,
counted with multiplicity, is the multiset \eqref{eq:spec} with multiplicities
\eqref{eq:multiplicity}.

Finally we check that \eqref{eq:spec} is well defined on the reduced source space, i.e.\ invariant
under $\theta\mapsto\theta+2\pi k$: the substitution $n\mapsto n+k$ gives
$\lambda_{n+k}(\theta+2\pi k)=\bigl(\tfrac{2\pi(n+k)-\theta-2\pi k}{L}\bigr)^2=\lambda_n(\theta)$,
and $n\mapsto n+k$ is a bijection of $\Zz$, so the multiset \eqref{eq:spec} is unchanged, with all
multiplicities preserved. Thus $\Spec(H_\theta)$ depends only on $\theta$ mod $2\pi$, i.e.\ only
on the point of $\Src_{\mathrm{flat}}$.
\end{proof}

\begin{remark}[Multiplicities are part of the answer]\label{rem:multiplicities}
We record \eqref{eq:spec} as a multiset, not as a subset of $\Rr$, because the multiplicity data
is physically real and is not recoverable from the underlying set. For generic $\theta$, distinct
$n$ give distinct $\lambda_n(\theta)$ and every eigenvalue is simple; but at $\theta=0$ the
eigenvalues $\lambda_n$ and $\lambda_{-n}$ coincide for all $n$, so every nonzero eigenvalue has
multiplicity two while $\lambda_0=0$ is simple; and at $\theta=\pi$ one finds
$\lambda_n(\pi)=\bigl(\tfrac{2\pi n-\pi}{L}\bigr)^2=\bigl(\tfrac{\pi(2n-1)}{L}\bigr)^2$, so
$\lambda_n(\pi)=\lambda_{1-n}(\pi)$ and every eigenvalue has multiplicity two. Discarding
multiplicities would erase this structure and, in particular, would change the target: ``the set
of eigenvalues'' and ``the spectrum with multiplicities'' are different correct-answer maps and
may have different completions. We take the latter throughout.
\end{remark}

\subsection{The spectral target and the curvature obstruction}

\begin{definition}[Spectral target]\label{def:spectral-target}
Let $A_{\mathrm{spec}}$ be the set of multisets of nonnegative reals arising as spectra of the
operators $H_\theta$, and define the correct-answer map
\begin{equation}\label{eq:aspec}
a_{\mathrm{spec}}\colon\Src_{\mathrm{flat}}\longrightarrow A_{\mathrm{spec}},
\qquad
a_{\mathrm{spec}}(\theta):=\Spec(H_\theta)
\end{equation}
with $\Spec(H_\theta)$ as in \eqref{eq:spec}. By the last paragraph of the proof of
\cref{prop:spectrum}, $a_{\mathrm{spec}}$ is well defined on $\Src_{\mathrm{flat}}$, i.e.\ on
$\theta$ mod $2\pi$. We denote this target by $Q_{\mathrm{spec}}$.
\end{definition}

\begin{theorem}[The spectral target is obstructed by the curvature
readout]\label{thm:curvature-obstruction}
$Q_{\mathrm{spec}}$ is not exactly solvable from the curvature readout $p_F$ of \eqref{eq:pF}.
Explicitly, the reduced source states $\theta=0$ and $\theta=\pi$ satisfy
\begin{equation}\label{eq:zero-pi}
p_F(0)=p_F(\pi)=0,
\qquad\text{but}\qquad
\Spec(H_0)\neq\Spec(H_\pi).
\end{equation}
Indeed $0\in\Spec(H_0)$ while $0\notin\Spec(H_\pi)$, and moreover
$\min\Spec(H_0)=0<\pi^2/L^2=\min\Spec(H_\pi)$.
\end{theorem}

\begin{proof}
First, $\theta=0$ and $\theta=\pi$ are distinct points of $\Src_{\mathrm{flat}}=\Rr/2\pi\Zz$,
since $\pi\not\equiv 0 \pmod{2\pi}$; equivalently their holonomies $h(0)=1$ and $h(\pi)=-1$ are
distinct elements of $\Uone$. They are therefore genuinely gauge-inequivalent physical
configurations, not two labels for one state.

Both are flat: by \eqref{eq:Ftheta}, $F_0=F_\pi=0$, so $p_F(0)=p_F(\pi)=0$ and the two states lie
in the same (indeed the only) fiber of $p_F$.

Now compute. By \eqref{eq:Dpsin} with $\theta=0$,
\begin{equation*}
\lambda_n(0)=\left(\frac{2\pi n}{L}\right)^{\!2},
\end{equation*}
and taking $n=0$ gives $\lambda_0(0)=0$, so $0\in\Spec(H_0)$; the corresponding eigenvector is
the constant function $\psi_0=L^{-1/2}$, and indeed $\Dx\psi_0=\partial_x\psi_0=0$ when
$\theta=0$.

By \eqref{eq:Dpsin} with $\theta=\pi$,
\begin{equation*}
\lambda_n(\pi)=\left(\frac{2\pi n-\pi}{L}\right)^{\!2}
=\frac{\pi^2}{L^2}\,(2n-1)^2 .
\end{equation*}
Since $2n-1$ is odd for every $n\in\Zz$, we have $(2n-1)^2\geq 1$, hence
\begin{equation}\label{eq:pi-lowerbound}
\lambda_n(\pi)\geq\frac{\pi^2}{L^2}>0
\qquad\text{for every } n\in\Zz,
\end{equation}
with equality precisely when $2n-1=\pm1$, i.e.\ $n\in\{0,1\}$. Therefore $0\notin\Spec(H_\pi)$
and $\min\Spec(H_\pi)=\pi^2/L^2$, attained with multiplicity two.

Since $0\in\Spec(H_0)$ and $0\notin\Spec(H_\pi)$, the two multisets differ, so
$a_{\mathrm{spec}}(0)\neq a_{\mathrm{spec}}(\pi)$ while $p_F(0)=p_F(\pi)$. This is exactly the
hypothesis \eqref{eq:obstruction-hyp} of \cref{cor:obstruction}, whence $Q_{\mathrm{spec}}$ is not
exactly solvable from $p_F$: no map $b\colon\{0\}\to A_{\mathrm{spec}}$ --- that is, no single
spectrum assigned to the curvature value $0$ --- is correct for both $\theta=0$ and $\theta=\pi$.
\end{proof}

\begin{corollary}[Holonomy is target-active for the spectrum]\label{cor:holonomy-active}
Let $r=h\colon\Src_{\mathrm{flat}}\to\Uone$, $h(\theta)=e^{i\theta}$, be the holonomy relation,
and let $\Phi_{\mathrm{spec}}\bigl(0,e^{i\theta}\bigr):=\Spec(H_\theta)$, which is well defined by
\cref{prop:spectrum}, so that $a_{\mathrm{spec}}=\Phi_{\mathrm{spec}}\circ(p_F,h)$ in the sense of
\eqref{eq:PhiQ}. Then $h$ is active for $Q_{\mathrm{spec}}$ on the single fiber $F_0$ of $p_F$.
\end{corollary}

\begin{proof}
Take $s_0=\theta=0$ and $s_1=\theta=\pi$ in \cref{def:active}. Both lie in $F_0$ by
\eqref{eq:pF-fiber}, and by \cref{thm:curvature-obstruction},
$\Phi_{\mathrm{spec}}(0,h(0))=\Spec(H_0)\neq\Spec(H_\pi)=\Phi_{\mathrm{spec}}(0,h(\pi))$, which is
\eqref{eq:active}.
\end{proof}

The contrast with \cref{ex:unobstructed} is the whole point. The readout $p_F$ is the same in both
cases, and the structural loss --- the collapse of all of $\Uone$ to a point --- is the same. For
the flatness target the loss is inactive and the readout is fully adequate. For the spectral
target the same loss is active, and by \cref{thm:active-obstructs} it obstructs. Whether
completion is necessary is a fact about the pair, never about the readout alone.

\begin{remark}[Physical reading]\label{rem:physical-reading}
The physical content of \cref{thm:curvature-obstruction} is the Aharonov--Bohm phenomenon in
spectral form \cite{AharonovBohm,WuYang}. Along the circle the field strength vanishes ---
identically, and for dimensional reasons --- so there is no local electromagnetic field anywhere
on the particle's accessible configuration space for it to feel. Nevertheless
the charged particle's energy levels depend on $\theta$: the covariant derivative
\eqref{eq:Dx} twists the periodicity condition felt by the wavefunction, shifting the Fourier
momenta from $2\pi n/L$ to $(2\pi n-\theta)/L$. At $\theta=\pi$ the shift is maximal, the
momentum lattice is displaced by half a spacing, no mode has zero covariant momentum, and the
ground-state energy is lifted from $0$ to $\pi^2/L^2$. The obstruction is not an artifact of the
formalism; it is a measurable spectral shift produced by information that curvature does not
record.
\end{remark}

%===============================================================
\section{The minimal completion for the spectral target}\label{sec:minimal-spectral}

Curvature is insufficient for $Q_{\mathrm{spec}}$. Holonomy is manifestly sufficient, since by
\cref{prop:spectrum} the spectrum is a function of $\theta$ mod $2\pi$. The question this section
answers is the sharp one: \emph{how much less than holonomy is enough?} The answer is: strictly
less. The spectral target sees holonomy only up to inversion.

\subsection{Holonomy is sufficient}

\begin{proposition}[Holonomy completes the curvature readout for the
spectrum]\label{prop:holonomy-sufficient}
The identity map $e=\id\colon\Src_{\mathrm{flat}}\to\Src_{\mathrm{flat}}$, equivalently the
holonomy map $h\colon\Src_{\mathrm{flat}}\to\Uone$, $h(\theta)=e^{i\theta}$, is a readout
extension of $p_F$ that is sufficient for $Q_{\mathrm{spec}}$.
\end{proposition}

\begin{proof}
$h$ is a bijection $\Src_{\mathrm{flat}}\to\Uone$ by \eqref{eq:Sflat} and \eqref{eq:holonomy},
hence surjective, and $p_F=u\circ h$ with $u\colon\Uone\to\{0\}$ the constant map, so $h$ is a
readout extension in the sense of \cref{def:extension}. Define
$b\colon\Uone\to A_{\mathrm{spec}}$ by $b(e^{i\theta}):=\Spec(H_\theta)$; this is well defined
because, as shown at the end of the proof of \cref{prop:spectrum}, $\Spec(H_\theta)$ depends only
on $\theta$ mod $2\pi$, i.e.\ only on $e^{i\theta}$. Then
$a_{\mathrm{spec}}=b\circ h$, so $h$ is $Q_{\mathrm{spec}}$-sufficient by \cref{def:sufficient}.
\end{proof}

\subsection{Spectral rigidity: the exact equivalence}

\begin{theorem}[Isospectrality iff $\theta'\equiv\pm\theta$]\label{thm:isospectral}
For $\theta,\theta'\in\Rr$,
\begin{equation}\label{eq:isospectral}
\Spec(H_\theta)=\Spec(H_{\theta'})
\qquad\Longleftrightarrow\qquad
\theta'\equiv\theta \ \ \text{or}\ \ \theta'\equiv-\theta \pmod{2\pi},
\end{equation}
where the spectra are compared as multisets, with multiplicities.
\end{theorem}

\begin{proof}
\emph{($\Leftarrow$) Reverse direction.} Suppose first $\theta'\equiv\theta\pmod{2\pi}$. Then
$\Spec(H_{\theta'})=\Spec(H_\theta)$ by the invariance established at the end of the proof of
\cref{prop:spectrum}.

Suppose next $\theta'\equiv-\theta \pmod{2\pi}$. By the same invariance it suffices to treat
$\theta'=-\theta$. Then for every $n\in\Zz$,
\begin{equation}\label{eq:reindex}
\lambda_{-n}(-\theta)
=\left(\frac{2\pi(-n)-(-\theta)}{L}\right)^{\!2}
=\left(\frac{-(2\pi n-\theta)}{L}\right)^{\!2}
=\left(\frac{2\pi n-\theta}{L}\right)^{\!2}
=\lambda_n(\theta).
\end{equation}
The map $n\mapsto -n$ is a bijection of $\Zz$ onto itself. Hence \eqref{eq:reindex} exhibits a
bijection between the indexing set of $\Spec(H_{-\theta})$ and that of $\Spec(H_\theta)$ carrying
each eigenvalue to an equal eigenvalue. A multiset is determined by its indexed family up to
bijective reindexing, so the two multisets are equal, \emph{with all multiplicities}:
\begin{equation*}
m_{-\theta}(\lambda)=\#\{n:\lambda_n(-\theta)=\lambda\}
=\#\{n:\lambda_{-n}(-\theta)=\lambda\}
=\#\{n:\lambda_{n}(\theta)=\lambda\}=m_\theta(\lambda).
\end{equation*}
So $\Spec(H_{-\theta})=\Spec(H_\theta)$. Concretely, the reindexing $n\mapsto -n$ is implemented
on the Hilbert space by complex conjugation, $\psi_n\mapsto\overline{\psi_n}=\psi_{-n}$, which is
an antiunitary intertwiner between $H_\theta$ and $H_{-\theta}$; but the elementary
reindexing argument above is already complete and requires nothing further.

\emph{($\Rightarrow$) Forward direction.} Suppose $\Spec(H_\theta)=\Spec(H_{\theta'})$. We
recover $\theta$ up to sign and $2\pi$ from the spectrum, using the lowest eigenvalue.

Define, for any $\theta\in\Rr$, the \emph{principal representative}
\begin{equation}\label{eq:alpha}
\alpha(\theta)\in[0,\pi]
\qquad\text{the unique number in } [0,\pi] \text{ with }
\theta\equiv\alpha(\theta) \ \text{ or }\ \theta\equiv-\alpha(\theta) \pmod{2\pi}.
\end{equation}
Existence and uniqueness of $\alpha(\theta)$: reduce $\theta$ mod $2\pi$ to a representative
$\tilde\theta\in[0,2\pi)$. If $\tilde\theta\in[0,\pi]$ set $\alpha:=\tilde\theta$; if
$\tilde\theta\in(\pi,2\pi)$ set $\alpha:=2\pi-\tilde\theta\in(0,\pi)$, and then
$\theta\equiv\tilde\theta=-\alpha \pmod{2\pi}$. In both cases $\alpha\in[0,\pi]$ and
$\theta\equiv\pm\alpha$. For uniqueness, suppose $\alpha,\beta\in[0,\pi]$ both work; then
$\alpha\equiv\pm\beta\pmod{2\pi}$. If $\alpha\equiv\beta$ then $\alpha-\beta\in 2\pi\Zz$ and
$|\alpha-\beta|\leq\pi<2\pi$, forcing $\alpha=\beta$. If $\alpha\equiv-\beta$ then
$\alpha+\beta\in2\pi\Zz$ with $0\le\alpha+\beta\le 2\pi$, so $\alpha+\beta\in\{0,2\pi\}$, giving
either $\alpha=\beta=0$ or $\alpha=\beta=\pi$. In all cases $\alpha=\beta$.

We claim
\begin{equation}\label{eq:lambdamin}
\lambda_{\min}(\theta):=\min\Spec(H_\theta)=\frac{\alpha(\theta)^2}{L^2}.
\end{equation}
To see this, note $\lambda_n(\theta)=\bigl(\tfrac{2\pi n-\theta}{L}\bigr)^2$, so
\begin{equation*}
\min_{n\in\Zz}\lambda_n(\theta)
=\frac{1}{L^2}\min_{n\in\Zz}\,(2\pi n-\theta)^2
=\frac{1}{L^2}\Bigl(\operatorname{dist}\bigl(\theta,\,2\pi\Zz\bigr)\Bigr)^{\!2},
\end{equation*}
the minimum being attained since $(2\pi n-\theta)^2\to\infty$ as $|n|\to\infty$. It remains to
identify $\operatorname{dist}(\theta,2\pi\Zz)=\alpha(\theta)$. Write $\theta\equiv\pm\alpha$ with
$\alpha=\alpha(\theta)\in[0,\pi]$; since $2\pi\Zz$ is invariant under negation and under
translation by $2\pi$, we get
$\operatorname{dist}(\theta,2\pi\Zz)=\operatorname{dist}(\pm\alpha,2\pi\Zz)
=\operatorname{dist}(\alpha,2\pi\Zz)$. For $\alpha\in[0,\pi]$ the nearest point of $2\pi\Zz$ is
$0$, since the next candidates are $\pm2\pi$ at distance $2\pi-\alpha\ge\pi\ge\alpha$ and
$2\pi+\alpha>\alpha$. Hence $\operatorname{dist}(\alpha,2\pi\Zz)=\alpha$, proving
\eqref{eq:lambdamin}.

Now assume $\Spec(H_\theta)=\Spec(H_{\theta'})$. Equal multisets have equal minima, so by
\eqref{eq:lambdamin}
\begin{equation*}
\frac{\alpha(\theta)^2}{L^2}=\lambda_{\min}(\theta)=\lambda_{\min}(\theta')
=\frac{\alpha(\theta')^2}{L^2},
\end{equation*}
whence $\alpha(\theta)^2=\alpha(\theta')^2$. Since $\alpha(\theta),\alpha(\theta')\in[0,\pi]$ are
both nonnegative, $\alpha(\theta)=\alpha(\theta')=:\alpha$. By \eqref{eq:alpha},
$\theta\equiv\pm\alpha$ and $\theta'\equiv\pm\alpha\pmod{2\pi}$. Combining the four sign cases:
if $\theta\equiv\alpha$ and $\theta'\equiv\alpha$ then $\theta'\equiv\theta$; if
$\theta\equiv\alpha$ and $\theta'\equiv-\alpha$ then $\theta'\equiv-\theta$; if
$\theta\equiv-\alpha$ and $\theta'\equiv\alpha$ then $\theta'\equiv-\theta$; if
$\theta\equiv-\alpha$ and $\theta'\equiv-\alpha$ then $\theta'\equiv\theta$. In every case
$\theta'\equiv\pm\theta\pmod{2\pi}$, as required.
\end{proof}

\begin{corollary}[The spectral target sees holonomy only modulo inversion]\label{cor:inversion}
Define the equivalence $\sim_{\mathrm{spec}}$ on $\Src_{\mathrm{flat}}\cong\Uone$ by
\begin{equation}\label{eq:spec-equiv}
e^{i\theta}\sim_{\mathrm{spec}}e^{i\theta'}
\qquad\Longleftrightarrow\qquad
\theta'\equiv\pm\theta\pmod{2\pi}
\qquad\Longleftrightarrow\qquad
e^{i\theta'}\in\{e^{i\theta},e^{-i\theta}\}.
\end{equation}
Then $\sim_{\mathrm{spec}}$ is exactly the answer equivalence of the spectral target:
$a_{\mathrm{spec}}(\theta)=a_{\mathrm{spec}}(\theta')$ iff $\theta\sim_{\mathrm{spec}}\theta'$.
\end{corollary}

\begin{proof}
Immediate from \cref{thm:isospectral} and \cref{def:spectral-target}.
\end{proof}

\subsection{The coarsest sufficient completion and its coordinate}

\begin{theorem}[$\cos\theta$ coordinatizes the minimal spectral
completion]\label{thm:cos-minimal}
Let
\begin{equation}\label{eq:c}
c\colon\Src_{\mathrm{flat}}\longrightarrow[-1,1],
\qquad
c(\theta):=\cos\theta .
\end{equation}
Then:
\begin{enumerate}
\item\label{it:c-welldef} $c$ is a well-defined surjection onto $[-1,1]$;
\item\label{it:c-classes} $c(\theta)=c(\theta')$ if and only if
$\theta\sim_{\mathrm{spec}}\theta'$; that is, $c$ realizes exactly the spectral equivalence
classes;
\item\label{it:c-suff} $c$ is a readout extension of $p_F$ that is sufficient for
$Q_{\mathrm{spec}}$;
\item\label{it:c-min} $c$ is a coarsest such extension: the induced map
$\varphi_{[-1,1]}\colon[-1,1]\to E_{Q_{\mathrm{spec}}}$ of \cref{thm:universal} is a bijection.
Consequently every $Q_{\mathrm{spec}}$-sufficient readout extension $e$ must separate source
states with distinct values of $\cos\theta$.
\end{enumerate}
\end{theorem}

\begin{proof}
\eqref{it:c-welldef} $\cos$ is $2\pi$-periodic, so $c$ descends to $\Rr/2\pi\Zz$; its image is
$[-1,1]$ since $\cos$ restricted to $[0,\pi]$ is already onto $[-1,1]$.

\eqref{it:c-classes} The function $\cos$ is strictly decreasing on $[0,\pi]$, hence injective
there. Given $\theta,\theta'$, write $\alpha=\alpha(\theta),\alpha'=\alpha(\theta')\in[0,\pi]$ as
in \eqref{eq:alpha}; since $\cos$ is even and $2\pi$-periodic, $\cos\theta=\cos\alpha$ and
$\cos\theta'=\cos\alpha'$. Therefore
\begin{equation*}
c(\theta)=c(\theta')
\iff \cos\alpha=\cos\alpha'
\iff \alpha=\alpha'
\iff \theta'\equiv\pm\theta \pmod{2\pi}
\iff \theta\sim_{\mathrm{spec}}\theta',
\end{equation*}
where the second equivalence uses injectivity of $\cos$ on $[0,\pi]$, and the third is the
argument at the end of the proof of \cref{thm:isospectral} (equality of principal representatives
is equivalent to $\theta'\equiv\pm\theta$).

\eqref{it:c-suff} $c$ is surjective by \eqref{it:c-welldef}, and $p_F=u\circ c$ with
$u\colon[-1,1]\to\{0\}$ constant, so $c$ is a readout extension. For sufficiency, define
$b_{\mathrm{spec}}\colon[-1,1]\to A_{\mathrm{spec}}$ by
\begin{equation}\label{eq:b-from-c}
b_{\mathrm{spec}}(\kappa):=\Spec\bigl(H_{\arccos\kappa}\bigr),
\qquad \arccos\kappa\in[0,\pi].
\end{equation}
If $c(\theta)=\kappa$ then by \eqref{it:c-classes} $\theta\sim_{\mathrm{spec}}\arccos\kappa$
(both have principal representative $\arccos\kappa$), so by \cref{cor:inversion}
$\Spec(H_\theta)=\Spec(H_{\arccos\kappa})=b_{\mathrm{spec}}(c(\theta))$. Hence
$a_{\mathrm{spec}}=b_{\mathrm{spec}}\circ c$, and $c$ is $Q_{\mathrm{spec}}$-sufficient.

\eqref{it:c-min} Write $\varphi:=\varphi_{[-1,1]}$ and $E_{\mathrm{spec}}
:=E_{Q_{\mathrm{spec}}}$ for brevity. By \cref{thm:universal}, applied to the sufficient
extension $c$, there is a unique surjection $\varphi\colon[-1,1]\to E_{\mathrm{spec}}$ with
$j_{Q_{\mathrm{spec}}}=\varphi\circ c$, given by
\begin{equation}\label{eq:varphi-c}
\varphi(\kappa)
=\bigl(u(\kappa),\,b_{\mathrm{spec}}(\kappa)\bigr)
=\bigl(0,\,b_{\mathrm{spec}}(\kappa)\bigr),
\qquad \kappa\in[-1,1].
\end{equation}
It remains to prove injectivity. Suppose $\varphi(\kappa)=\varphi(\kappa')$. Comparing second
coordinates in \eqref{eq:varphi-c} gives $b_{\mathrm{spec}}(\kappa)=b_{\mathrm{spec}}(\kappa')$,
which by \eqref{eq:b-from-c} says
\begin{equation}\label{eq:iso-arccos}
\Spec\bigl(H_{\arccos\kappa}\bigr)=\Spec\bigl(H_{\arccos\kappa'}\bigr).
\end{equation}
By \cref{thm:isospectral}, \eqref{eq:iso-arccos} forces
$\arccos\kappa'\equiv\pm\arccos\kappa \pmod{2\pi}$. Both $\arccos\kappa$ and $\arccos\kappa'$
lie in $[0,\pi]$, so each is its own principal representative in the sense of \eqref{eq:alpha};
by the uniqueness established there, they are equal. Applying $\cos$ yields $\kappa=\kappa'$.

Hence $\varphi$ is a bijection, and $c$ carries precisely the distinctions of
$E_{\mathrm{spec}}$ --- no more and no fewer.

The final assertion follows from \cref{cor:ordering}: if $e$ is $Q_{\mathrm{spec}}$-sufficient and
$c(\theta_0)\neq c(\theta_1)$, then by \eqref{it:c-classes} $\theta_0\not\sim_{\mathrm{spec}}
\theta_1$, so by \cref{cor:inversion} $a_{\mathrm{spec}}(\theta_0)\neq a_{\mathrm{spec}}
(\theta_1)$, and \cref{thm:separation} forces $e(\theta_0)\neq e(\theta_1)$.
\end{proof}

\begin{remark}[What is and is not claimed about $\cos\theta$]\label{rem:cos-caveat}
\Cref{thm:cos-minimal} says that $\cos\theta$ \emph{coordinatizes} the coarsest completion
sufficient for $Q_{\mathrm{spec}}$: its level sets are exactly the spectral equivalence classes,
and it separates exactly what must be separated. It does not say that $\cos\theta$ is the unique
possible numerical coordinate on $E_{Q_{\mathrm{spec}}}$. Any injective function of $\cos\theta$
--- for instance $\tfrac12(1-\cos\theta)$, or $\alpha(\theta)$ itself, or the ground-state energy
$\lambda_{\min}=\alpha(\theta)^2/L^2$ --- has the same level sets and coordinatizes the same
quotient equally well; the completion is the quotient, and the coordinate is a labeling of it.
Nor does \cref{thm:cos-minimal} say that $\cos\theta$ is minimal for \emph{every} physical target
on $\Src_{\mathrm{flat}}$. It is minimal for the nominated spectral target, and for that target
only.
\end{remark}

\begin{example}[A target for which $\cos\theta$ is insufficient: oriented
transport]\label{ex:transport}
Consider the target $Q_{\Hol}$ with $A_{Q_{\Hol}}=\Uone$ and correct-answer map
$a_{\Hol}(\theta)=e^{i\theta}$: ``what phase does a unit charge acquire on being transported once
around the circle in the positive direction?'' Take $\theta=\pi/2$ and $\theta'=-\pi/2$. Then
\begin{equation*}
c(\pi/2)=\cos(\pi/2)=0=\cos(-\pi/2)=c(-\pi/2),
\end{equation*}
so $c$ does not separate these two states; but
$a_{\Hol}(\pi/2)=e^{i\pi/2}=i\neq -i=e^{-i\pi/2}=a_{\Hol}(-\pi/2)$. By \cref{thm:separation}, $c$
is not $Q_{\Hol}$-sufficient. By \cref{thm:isospectral} the two states are nonetheless
isospectral, $\Spec(H_{\pi/2})=\Spec(H_{-\pi/2})$, consistent with $c$ being sufficient for
$Q_{\mathrm{spec}}$.

The full holonomy map $h$ \emph{is} $Q_{\Hol}$-sufficient (take $b=\id_{\Uone}$), and by
\cref{cor:ordering} it is coarsest for that target, since $j_{Q_{\Hol}}=(p_F,a_{\Hol})$ has the
same level sets as $h$. So on one and the same reduced source space, with one and the same
readout, the minimal completion for the transport target is the full holonomy circle, while the
minimal completion for the spectral target is its strictly coarser inversion quotient
$\Uone/(z\sim\bar z)\cong[-1,1]$. The scalar spectrum is blind to orientation; oriented transport
is not. This is target-relativity of completion in its sharpest form, and it is not a defect of
either description: each retains exactly what its question requires.
\end{example}

%===============================================================
\section{Finite-time dynamical descent}\label{sec:dynamics}

Our final application of the criterion concerns the question that motivates most practical
coarse-graining: can the projected description evolve on its own?

\subsection{Setting}

Let $\Phi_\tau\colon\Src\to\Src$ denote the source evolution through elapsed time $\tau$, assumed
to be a well-defined map for the fixed $\tau$ under consideration. No group law, continuity, or
reversibility is needed for what follows; $\tau$ is fixed throughout this section and we prove a
finite-time statement. Let $p\colon\Src\to\Rdt$ be the readout as before.

\begin{definition}[Future-readout target]\label{def:future-target}
The \emph{future-readout target} at time $\tau$ is the target $Q_\tau$ with answer set
$A_{Q_\tau}=\Rdt$ and correct-answer map
\begin{equation}\label{eq:atau}
a_\tau:=p\circ\Phi_\tau\colon\Src\longrightarrow\Rdt,
\qquad
a_\tau(s)=p\bigl(\Phi_\tau(s)\bigr).
\end{equation}
In words: ``given the source state now, what will the readout be after time $\tau$?''
\end{definition}

\begin{definition}[Autonomous readout propagator]\label{def:propagator}
An \emph{autonomous readout propagator} at time $\tau$ is a map $U_\tau\colon\Rdt\to\Rdt$ such
that the projection intertwines the dynamics:
\begin{equation}\label{eq:intertwine}
p\circ\Phi_\tau=U_\tau\circ p .
\end{equation}
\end{definition}

Existence of $U_\tau$ is precisely the statement that the readout description is dynamically
closed at time $\tau$: the future readout value can be computed from the present readout value
alone, without consulting the source.

\subsection{The closure criterion}

\begin{theorem}[Existence criterion for the autonomous propagator]\label{thm:closure}
An autonomous readout propagator $U_\tau$ satisfying \eqref{eq:intertwine} exists if and only if
\begin{equation}\label{eq:closure-condition}
p(s)=p(s')
\qquad\Longrightarrow\qquad
p\bigl(\Phi_\tau(s)\bigr)=p\bigl(\Phi_\tau(s')\bigr)
\qquad\text{for all } s,s'\in\Src .
\end{equation}
When it exists, $U_\tau$ is unique.
\end{theorem}

\begin{proof}
Equation \eqref{eq:intertwine} states exactly that $a_\tau=p\circ\Phi_\tau$ factors through $p$,
i.e.\ that the target $Q_\tau$ of \cref{def:future-target} is exactly solvable from the readout
with $\bar a_{Q_\tau}=U_\tau$. By \cref{thm:exact-solvability} applied to $Q_\tau$ (with
$A_{Q_\tau}=\Rdt$ and $a_{Q_\tau}=a_\tau$), such a factorization exists if and only if $a_\tau$
is constant on every fiber of $p$, which is condition \eqref{it:imp} of that theorem and reads,
after substituting $a_\tau=p\circ\Phi_\tau$, exactly as \eqref{eq:closure-condition}. Uniqueness
of $U_\tau$ is the uniqueness clause of \cref{thm:exact-solvability}.
\end{proof}

\begin{corollary}[Dynamical obstruction]\label{cor:dynamical-obstruction}
Suppose there exist $s_0,s_1\in\Src$ with
\begin{equation}\label{eq:dyn-obstruction}
p(s_0)=p(s_1)
\qquad\text{but}\qquad
p\bigl(\Phi_\tau(s_0)\bigr)\neq p\bigl(\Phi_\tau(s_1)\bigr).
\end{equation}
Then no map $U_\tau\colon\Rdt\to\Rdt$ satisfies \eqref{eq:intertwine}. Concretely: two source
states with the \emph{same present readout} but \emph{different future readouts} cannot both be
reproduced by any deterministic evolution law acting on the present readout state alone.
\end{corollary}

\begin{proof}
Apply \cref{cor:obstruction} to the target $Q_\tau$: \eqref{eq:dyn-obstruction} is precisely
\eqref{eq:obstruction-hyp} for $a_\tau$. For any candidate $U_\tau$, the value
$U_\tau(t)$ at the common readout $t=p(s_0)=p(s_1)$ is a single element of $\Rdt$ and hence
cannot equal both $p(\Phi_\tau(s_0))$ and $p(\Phi_\tau(s_1))$.
\end{proof}

\subsection{Prediction at time \texorpdfstring{$\tau$}{tau} versus autonomous evolution}
\label{subsec:predict-vs-evolve}

The general theory of \cref{sec:separation,sec:joint} applies to $Q_\tau$ verbatim, and it is
worth stating exactly what it delivers --- and, just as importantly, what it does not.

By \cref{thm:separation}, any readout extension $e\colon\Src\to E$ sufficient for $Q_\tau$ must
separate the obstructing pair $s_0,s_1$ of \eqref{eq:dyn-obstruction}. By
\cref{prop:EQ-sufficient,thm:universal}, the coarsest such extension is the joint target image
\begin{equation}\label{eq:EQtau}
E_{Q_\tau}=\Img\bigl(p,\;p\circ\Phi_\tau\bigr)\subseteq\Rdt\times\Rdt,
\end{equation}
the present readout paired with the correct future readout at time $\tau$.

\begin{proposition}[What $Q_\tau$-sufficiency provides]\label{prop:qtau-provides}
Let $e\colon\Src\to E$ be a readout extension sufficient for $Q_\tau$. Then there is a unique map
\begin{equation}\label{eq:btau}
b_\tau\colon E\longrightarrow\Rdt
\qquad\text{with}\qquad
p\circ\Phi_\tau=b_\tau\circ e ,
\end{equation}
so that the $E$-value at the present instant determines the readout value at time $\tau$.
\end{proposition}

\begin{proof}
This is \cref{def:sufficient} and \cref{lem:sufficiency-criterion} applied to $Q_\tau$, whose
answer set is $A_{Q_\tau}=\Rdt$ and whose correct-answer map is $a_\tau=p\circ\Phi_\tau$.
\end{proof}

\Cref{prop:qtau-provides} is a statement about \emph{prediction of a $\Rdt$-valued quantity at the
one specified time $\tau$}. It says the enriched present state determines the future readout. It
does \emph{not} say that the enriched state evolves --- that is, it produces a map $E\to\Rdt$, not
a map $E\to E$, and there is no sense in which \eqref{eq:btau} can be iterated. Autonomous
evolution on $E$ is a strictly stronger requirement, and it is governed by its own descent
condition.

\begin{definition}[Autonomous evolution on an extension]\label{def:autonomous-E}
Let $e\colon\Src\to E$ be a readout extension. An \emph{autonomous evolution on $E$ at time
$\tau$} is a map $V_\tau\colon E\to E$ satisfying
\begin{equation}\label{eq:E-descent}
e\circ\Phi_\tau=V_\tau\circ e .
\end{equation}
\end{definition}

\begin{proposition}[Criterion for autonomous evolution on $E$]\label{prop:autonomous-criterion}
A map $V_\tau$ satisfying \eqref{eq:E-descent} exists if and only if
\begin{equation}\label{eq:E-descent-condition}
e(s)=e(s')
\qquad\Longrightarrow\qquad
e\bigl(\Phi_\tau(s)\bigr)=e\bigl(\Phi_\tau(s')\bigr)
\qquad\text{for all } s,s'\in\Src ,
\end{equation}
and it is then unique.
\end{proposition}

\begin{proof}
Apply \cref{thm:exact-solvability} to the surjection $e$ and the map $e\circ\Phi_\tau\colon
\Src\to E$: a factorization $e\circ\Phi_\tau=V_\tau\circ e$ exists iff $e\circ\Phi_\tau$ is
constant on the fibers of $e$, which is \eqref{eq:E-descent-condition}, and $V_\tau$ is then
unique by the uniqueness clause of that theorem.
\end{proof}

Conditions \eqref{eq:btau} and \eqref{eq:E-descent} are different requirements, and the
implication between them runs in exactly one direction.

\begin{proposition}[Autonomous evolution implies $Q_\tau$-sufficiency]\label{prop:auto-implies-suff}
Let $e\colon\Src\to E$ be a readout extension, with structure map $u\colon E\to\Rdt$ satisfying
$p=u\circ e$. If $e$ admits an autonomous evolution $V_\tau\colon E\to E$ at time $\tau$, in the
sense of \eqref{eq:E-descent}, then $e$ is sufficient for $Q_\tau$, and one may take
\begin{equation}\label{eq:btau-from-V}
b_\tau=u\circ V_\tau .
\end{equation}
\end{proposition}

\begin{proof}
Using $p=u\circ e$ and then \eqref{eq:E-descent},
\begin{equation*}
p\circ\Phi_\tau
=u\circ e\circ\Phi_\tau
=u\circ V_\tau\circ e
=\bigl(u\circ V_\tau\bigr)\circ e ,
\end{equation*}
which is \eqref{eq:btau} with $b_\tau:=u\circ V_\tau\colon E\to\Rdt$. Hence $a_\tau=p\circ\Phi_\tau$
factors through $e$, and $e$ is $Q_\tau$-sufficient by \cref{def:sufficient}.
\end{proof}

The converse fails in general. Sufficiency for $Q_\tau$ constrains $e$ to determine the future
\emph{readout} --- a $\Rdt$-valued quantity --- whereas \eqref{eq:E-descent-condition} constrains
$e$ to determine the future \emph{$E$-value}, which is in general a finer quantity, since $e$
refines $p$ (\cref{lem:extension-basic}\,\eqref{it:refine}). Knowing where the present $E$-state
sends the time-$\tau$ readout says nothing about where it sends the time-$\tau$ $E$-state. The
following finite example shows that the gap is real and not merely unproved.

\begin{example}[A $Q_\tau$-sufficient extension admitting no autonomous evolution]
\label{ex:suff-not-autonomous}
Let
\begin{equation*}
\Src=\{a,b,c,d\},
\qquad
\Rdt=\{*\},
\end{equation*}
with $p\colon\Src\to\Rdt$ the constant readout, and let
\begin{equation*}
E=\{x,y,z\},
\qquad
e(a)=e(b)=x,
\quad
e(c)=y,
\quad
e(d)=z .
\end{equation*}
Then $e$ is surjective and $p=u\circ e$ for the constant map $u\colon E\to\Rdt$, so $e$ is a
readout extension. Define a source evolution $\Phi_\tau\colon\Src\to\Src$ by
\begin{equation*}
\Phi_\tau(a)=c,
\qquad
\Phi_\tau(b)=d,
\end{equation*}
with $\Phi_\tau(c)$ and $\Phi_\tau(d)$ chosen arbitrarily --- say $\Phi_\tau(c)=c$ and
$\Phi_\tau(d)=d$.

Since $p$ is constant, the composite $p\circ\Phi_\tau$ is constant, hence constant on every fiber
of $e$; by \cref{lem:sufficiency-criterion} it factors through $e$, so $e$ is $Q_\tau$-sufficient.
(Concretely, $b_\tau\colon E\to\Rdt$ is the constant map to $*$.)

Nevertheless $e$ admits no autonomous evolution. We have $e(a)=e(b)=x$, so $a$ and $b$ lie in a
common fiber of $e$, while
\begin{equation*}
e\bigl(\Phi_\tau(a)\bigr)=e(c)=y
\;\neq\;
z=e(d)=e\bigl(\Phi_\tau(b)\bigr).
\end{equation*}
Thus $e\circ\Phi_\tau$ is \emph{not} constant on the fiber $e^{-1}(x)$, so
\eqref{eq:E-descent-condition} fails, and by \cref{prop:autonomous-criterion} no map
$V_\tau\colon E\to E$ satisfies $e\circ\Phi_\tau=V_\tau\circ e$. Directly: such a $V_\tau$ would
have to send the single element $x$ to both $y$ and $z$.
\end{example}

So $Q_\tau$-sufficiency does not imply autonomous $E$-evolution, and autonomy is strictly stronger
in general, with $Q_\tau$-sufficiency the weaker consequence of it established in
\cref{prop:auto-implies-suff}.

\begin{remark}[What $E_{Q_\tau}$ is, and what it is not]\label{rem:EQtau-scope}
Accordingly, the joint target image \eqref{eq:EQtau} should be described in exactly one way: it is
the universally coarsest completion of the readout \emph{for predicting the $\Rdt$-valued readout
at the fixed time $\tau$}. By \cref{thm:universal} every $Q_\tau$-sufficient extension maps onto
it, and by \cref{prop:EQ-sufficient} it is itself sufficient for that prediction.

Its $Q_\tau$-sufficiency does not by itself make it an autonomous state space, supply an iterable
dynamics, or furnish a semigroup completion. Nothing proved here supplies a map
$E_{Q_\tau}\to E_{Q_\tau}$ satisfying \eqref{eq:E-descent}, and the results of this section do not
assert one. By \cref{prop:auto-implies-suff}, autonomy would be a strictly stronger property than
the $Q_\tau$-sufficiency that \cref{thm:universal} establishes for $E_{Q_\tau}$, and
\cref{ex:suff-not-autonomous} shows that the two can come apart; the stronger property is not
established here.

This leaves the question open rather than settled against it. In a particular system,
$E_{Q_\tau}$ may or may not happen to admit an autonomous $V_\tau$; what does not hold is that any
such evolution \emph{follows} from the universal construction of \cref{thm:universal}. Whether a
given extension admits an autonomous $V_\tau$, whether such extensions can be chosen compatibly
across a range of times $\tau$, and whether the resulting family composes as a semigroup are all
questions beyond \cref{thm:closure,prop:autonomous-criterion}, which are finite-time factorization
statements at a single fixed $\tau$. Those questions belong to \cite{ShadowV} and are not
addressed here.
\end{remark}

\subsection{What the obstruction does and does not license}

\Cref{cor:dynamical-obstruction} rules out one thing precisely: a deterministic autonomous law on
$\Rdt$ alone. It is silent about what to do instead, and it is worth being explicit about the
available modeling responses, none of which is an automatic consequence of the theorem.

\begin{itemize}
\item \emph{State refinement.} One may replace $\Rdt$ by a $Q_\tau$-sufficient extension $E$, on
which, by \cref{prop:qtau-provides}, a deterministic \emph{prediction} $b_\tau\colon E\to\Rdt$ of
the time-$\tau$ readout does exist. Obtaining a deterministic \emph{evolution} on $E$ is a further
demand, requiring \eqref{eq:E-descent-condition}, and is not delivered by $Q_\tau$-sufficiency.
Whether an extension meeting either condition is available in practice, and whether its variables
are measurable, are separate questions the theorem does not settle. What the theorem does settle
is a lower bound on what $E$ must distinguish.

\item \emph{Memory.} One may allow the evolution law to depend on the past history of the readout
rather than on its instantaneous value. This changes the shape of the modeling problem --- the
state is now a trajectory segment, not a point of $\Rdt$ --- and whether the history supplies the
missing distinctions is not settled by \cref{thm:closure}, which concerns the instantaneous
readout. Projection-operator methods in nonequilibrium statistical mechanics
\cite{Zwanzig,Grabert} are the classical setting in which such history-dependent reduced
equations arise; the construction of exact memory equations for the present framework belongs to
\cite{ShadowV} and is not developed here.

\item \emph{Probabilistic modeling.} One may posit a probability law and predict a distribution
over future readouts. See \cref{rem:stochastic}.

\item \emph{Domain restriction.} One may restrict $\Src$ to a subset on which
\eqref{eq:closure-condition} does hold, obtaining an exactly closed law on the restricted
problem. This is legitimate and often useful, but it is a change of the source domain, and the
theorem then applies to the restricted problem rather than the original one.
\end{itemize}

\begin{remark}[Stochasticity is not exact trajectory recovery]\label{rem:stochastic}
It is sometimes suggested that the obstruction of \cref{cor:dynamical-obstruction} is dissolved by
making the reduced law stochastic. It is not, and the reason is the one already given in
\cref{rem:not-probability}. Suppose $p(s_0)=p(s_1)=t$ but the futures differ. A stochastic law on
$\Rdt$ assigns to $t$ a distribution over $\Rdt$. If that law is intended not to exclude any
source-compatible future --- and a law that did exclude one would be predicting the impossibility
of something that can in fact occur --- then its support must contain both $p(\Phi_\tau(s_0))$ and
$p(\Phi_\tau(s_1))$. It therefore does not tell us which of the two \emph{actually occurs}, and if
the system is in fact in state $s_0$, the stochastic law has not predicted the realized
trajectory; it has predicted a spread that contains it. This is a strictly weaker claim, and it
answers a different target --- ``what is the distribution of the future readout, given the present
readout \emph{and a specified law on the fibers}?'' --- than the one we posed.

Note also the italicized clause. Even the distributional target requires input the readout does not
contain: a probability law on the fibers $F_t$, which is a piece of source-level information. Absent
such a law, the most that can be said from the readout alone is the set-valued statement of
\cref{thm:minimal-sound}, namely that the future readout lies in
$\CompatSet_{Q_\tau}(t)=p\bigl(\Phi_\tau(F_t)\bigr)$, and that this set is the least sound
prediction available. Stochastic modeling is a legitimate and powerful response to
\cref{cor:dynamical-obstruction}; it is not a refutation of it, and it should not be presented as
recovering the realized state.
\end{remark}

%===============================================================
\section{Discussion and scope}\label{sec:discussion}

\subsection{What has been established}

The results above are elementary in their machinery and, we would argue, correspondingly robust.
Their content can be compressed into four statements.

\emph{Obstruction is target-relative.} A noninjective readout is not, as such, defective. It fails
a target exactly when the correct answer to that target varies within one of its fibers
(\cref{thm:exact-solvability,cor:obstruction}) --- that is the general criterion, and it is stated
without auxiliary hypotheses. When a mediating variable $r$ is nominated, so that
$a_Q=\Phi_Q(p,r)$ as in \eqref{eq:PhiQ}, the same criterion localizes: the loss of $r$ obstructs
$Q$ exactly when $r$ is active for $Q$ on some fiber
(\cref{thm:active-obstructs,thm:inactive-preserves}). The same readout may pass one target and
fail another (\cref{ex:unobstructed} versus \cref{thm:curvature-obstruction}). Consequently the
inference from ``information was lost'' to ``the description must be completed'' is invalid
without a nominated target, and we have been careful never to make it.

\emph{Impossibility is derived, not assumed.} Every impossibility claim in this paper reduces to
the same counting observation: a deterministic construction from the readout has the form
$b\circ p$ and therefore assigns one answer per fiber, while the correct-answer map may assign
two. No hypothesis of the form ``no replacement exists'' is anywhere invoked, and
\cref{prop:auxiliary-useless} closes the natural attempt to evade the conclusion by manufacturing
derived readout variables.

\emph{Without completion, there is still a best answer, and it is a set.} The compatible-answer
set $\CompatSet_Q(t)=a_Q(F_t)$ is the pointwise-least uniformly sound set-valued rule
(\cref{thm:minimal-sound}), and exactness is equivalent to its being everywhere a singleton
(\cref{cor:singleton}). This is a set-theoretic fact prior to any probabilistic modeling; a
probability law over $\CompatSet_Q(t)$ ranks its elements but does not merge them and does not
identify which source state obtained (\cref{rem:not-probability,rem:stochastic}).

\emph{With completion, there is a least sufficient one.} Every $Q$-sufficient extension must
separate states with different answers (\cref{thm:separation}), and the joint target image
$E_Q=\Img(p,a_Q)$ both is sufficient (\cref{prop:EQ-sufficient}) and receives a unique surjection
from every sufficient extension (\cref{thm:universal}). It is thus the coarsest completion for the
target, and it retains the readout value and the answer distinction and nothing else. In
particular it does not reconstruct the source: \cref{cor:ordering} runs $E\to E_Q$, and the
identity extension $\id_\Src$ is one of the $E$'s.

\subsection{The example as a test of the theory}

The flat $\Uone$ example is not decoration. It exhibits, in a system where every object is
computable in closed form, all four statements at once. The reduced source space
$\Src_{\mathrm{flat}}\cong\Uone$ is a genuine moduli space, obtained only after quotienting by
large gauge transformations \eqref{eq:largegauge-action}, so its points are physically
inequivalent. The curvature readout \eqref{eq:pF} collapses it entirely. Against the flatness
target the collapse is inert (\cref{ex:unobstructed}); against the spectral target it is fatal
(\cref{thm:curvature-obstruction}), and the fatality is a measurable one --- the ground-state
energy jumps from $0$ to $\pi^2/L^2$ between $\theta=0$ and $\theta=\pi$, with the eigenvalue
multiplicities changing as recorded in \cref{rem:multiplicities}. Yet the completion demanded by
the spectral target is strictly less than the source: \cref{thm:isospectral} shows the spectrum
determines $\theta$ only up to sign, so \cref{thm:cos-minimal} identifies the coarsest sufficient
completion as the inversion quotient, coordinatized by $\cos\theta$. And this coarsest completion
is genuinely insufficient for a different, entirely physical target on the same space
(\cref{ex:transport}): oriented transport requires the whole holonomy circle, because it can tell
$e^{i\pi/2}$ from $e^{-i\pi/2}$ and the scalar spectrum cannot.

\subsection{Scope, and what is not claimed}

We close by delimiting the results, since each of the following overreaches is easy and none is
supported.

We have not shown that every lossy description requires repair; \cref{thm:inactive-preserves}
shows the opposite is routinely the case. We have not shown that the source is recoverable from
any completion; \cref{cor:ordering} explicitly denies it. We have not shown that curvature is a
poor physical observable, only that it is a poor readout \emph{for the spectral question}, which
is a claim about a pair (\cref{rem:curvature-restricted}). We have not shown that $\cos\theta$ is
canonical among numerical coordinates or minimal for all targets
(\cref{rem:cos-caveat,ex:transport}). We have not shown that stochastic reduced dynamics is
illegitimate, only that it does not recover the realized trajectory and does not answer the target
we posed (\cref{rem:stochastic}). And we have not made, and do not make, any claim to have
established that physical reality is a projection of a hidden source; the mathematics here
constrains what a given projection can and cannot deliver, given a nominated target, and that is
all it does.

Within the sequence of which this paper is a part, the role of the present results is to supply
the discrimination that the descent theorem of \cite{ShadowI} alone does not provide. Paper 1
identified when a source relation descends; this paper identifies when the failure to descend
matters, and what a repair must contain. Two natural generalizations are deliberately deferred.
Passing from a single correct-answer map to a family of realization relations, and constructing the
canonical minimal completion for such a family, requires machinery not developed here and is the
subject of \cite{ShadowIII}; the results above concern a single target throughout, and their
minimality claims should be read accordingly. Passing from the finite-time statement of
\cref{thm:closure} to continuous-time or semigroup dynamics, and to the exact memory equations
that govern projected evolution when \eqref{eq:closure-condition} fails, likewise lies outside the
present scope \cite{ShadowV}; \cref{sec:dynamics} proves exactly one thing about dynamics --- the
existence criterion for the finite-time propagator $U_\tau$ --- and no more.

%===============================================================
\section{Conclusion}\label{sec:conclusion}

The information a readout discards is not, by itself, an obstruction. It becomes one against a
nominated target, and it does so under an exactly identifiable condition: the correct answer must
vary across a single readout fiber. When that condition fails, the coarse description is complete
for the question at hand, however much it has forgotten
(\cref{thm:inactive-preserves,ex:unobstructed}). When it holds, no deterministic construction from
the readout --- and no auxiliary quantity computed from the readout --- can be correct throughout
(\cref{cor:obstruction,prop:auxiliary-useless}), and the honest residual is the compatible-answer
set, which is the least sound set-valued rule available and is not a probability
(\cref{thm:minimal-sound}).

Repair, when required, is then constrained from below, and the lower bound is attained
abstractly within the paper's category of extensions. Any sufficient enrichment must separate
every pair of states whose correct answers differ (\cref{thm:separation}), and hence must
distinguish at least the joint $(p,a_Q)$-classes (\cref{cor:ordering}). The joint target image
$E_Q=\Img(p,a_Q)$ attains this lower bound exactly as a set-theoretic construction: it is
sufficient for $Q$ (\cref{prop:EQ-sufficient}), it retains no information beyond the original
readout value and the target answer, and every sufficient enrichment maps onto it
(\cref{thm:universal}). This characterizes the least distinctions any repair must carry; it does
not by itself construct a physically measurable repair, which must independently supply the
enriched extension and the answer map that factors $a_Q$. The flat-connection example below does
supply one concretely.

This is a statement about the floor, not the ceiling. Other $Q$-sufficient extensions may --- and
in practice usually do --- carry additional information beyond the $(p,a_Q)$-classes: the identity
extension $\id_\Src$ retains the entire source, and any extension adequate for several targets at
once will generally be strictly finer than $E_Q$. Nothing forbids this, and \cref{thm:universal}
does not assert otherwise; it asserts only that whatever else such an extension carries, it maps
canonically onto $E_Q$, which is the least it could have carried and still answered $Q$.

The flat $\Uone$ connections on a circle make all of this concrete and, we think, persuasive. A
curvature readout that erases the entire flat moduli space is perfectly adequate for some
questions and provably inadequate for the spectrum of the charged covariant Laplacian, where
$H_0$ has a zero mode and $H_\pi$ has none. The completion the spectrum demands is not the source
and is not even the full holonomy: the spectrum determines the holonomy only up to inversion, and
the coarsest sufficient completion is $\Uone/(z\sim\bar z)$, coordinatized by $\cos\theta$ ---
strictly coarser than the holonomy that an oriented transport question would require. Two targets,
one source, one readout, two different minimal completions. That is the shape of the phenomenon,
and the same criterion, applied to the future readout, decides when a projected dynamics can close
on itself.

%===============================================================
\begin{thebibliography}{99}

\bibitem{AharonovBohm}
Y.~Aharonov and D.~Bohm,
\emph{Significance of electromagnetic potentials in the quantum theory},
Physical Review \textbf{115} (1959), 485--491.

\bibitem{BaezMuniain}
J.~Baez and J.~P.~Muniain,
\emph{Gauge Fields, Knots and Gravity},
World Scientific, Singapore, 1994.

\bibitem{Kobayashi}
S.~Kobayashi and K.~Nomizu,
\emph{Foundations of Differential Geometry, Volume I},
Interscience Publishers, New York, 1963.

\bibitem{NakaharaBook}
M.~Nakahara,
\emph{Geometry, Topology and Physics},
2nd ed., Institute of Physics Publishing, Bristol, 2003.

\bibitem{Grabert}
H.~Grabert,
\emph{Projection Operator Techniques in Nonequilibrium Statistical Mechanics},
Springer, Berlin, 1982.

\bibitem{ReedSimonI}
M.~Reed and B.~Simon,
\emph{Methods of Modern Mathematical Physics I: Functional Analysis},
revised and enlarged ed., Academic Press, New York, 1980.

\bibitem{ReedSimonIV}
M.~Reed and B.~Simon,
\emph{Methods of Modern Mathematical Physics IV: Analysis of Operators},
Academic Press, New York, 1978.

\bibitem{ShadowI}
J.~Rodgers,
\emph{Source--Readout Non-Equivalence: Descent and Equivariant Reconstruction Obstructions},
Zenodo, 2026.
\href{https://doi.org/10.5281/zenodo.21369967}{doi:10.5281/zenodo.21369967}.

\bibitem{ShadowIII}
J.~Rodgers,
\emph{Canonical minimal completion for families of realization relations},
forthcoming manuscript, 2026.

\bibitem{ShadowV}
J.~Rodgers,
\emph{Projected dynamics, semigroups, and exact memory equations},
forthcoming manuscript, 2026.

\bibitem{WuYang}
T.~T.~Wu and C.~N.~Yang,
\emph{Concept of nonintegrable phase factors and global formulation of gauge fields},
Physical Review D \textbf{12} (1975), 3845--3857.

\bibitem{Zwanzig}
R.~Zwanzig,
\emph{Ensemble method in the theory of irreversibility},
Journal of Chemical Physics \textbf{33} (1960), 1338--1341.

\end{thebibliography}

\end{document}
